\documentclass[11pt]{article}

    \usepackage[breakable]{tcolorbox}
    \usepackage{parskip} % Stop auto-indenting (to mimic markdown behaviour)
    

    % Basic figure setup, for now with no caption control since it's done
    % automatically by Pandoc (which extracts ![](path) syntax from Markdown).
    \usepackage{graphicx}
    % Keep aspect ratio if custom image width or height is specified
    \setkeys{Gin}{keepaspectratio}
    % Maintain compatibility with old templates. Remove in nbconvert 6.0
    \let\Oldincludegraphics\includegraphics
    % Ensure that by default, figures have no caption (until we provide a
    % proper Figure object with a Caption API and a way to capture that
    % in the conversion process - todo).
    \usepackage{caption}
    \DeclareCaptionFormat{nocaption}{}
    \captionsetup{format=nocaption,aboveskip=0pt,belowskip=0pt}

    \usepackage{float}
    \floatplacement{figure}{H} % forces figures to be placed at the correct location
    \usepackage{xcolor} % Allow colors to be defined
    \usepackage{enumerate} % Needed for markdown enumerations to work
    \usepackage{geometry} % Used to adjust the document margins
    \usepackage{amsmath} % Equations
    \usepackage{amssymb} % Equations
    \usepackage{textcomp} % defines textquotesingle
    % Hack from http://tex.stackexchange.com/a/47451/13684:
    \AtBeginDocument{%
        \def\PYZsq{\textquotesingle}% Upright quotes in Pygmentized code
    }
    \usepackage{upquote} % Upright quotes for verbatim code
    \usepackage{eurosym} % defines \euro

    \usepackage{iftex}
    \ifPDFTeX
        \usepackage[T1]{fontenc}
        \IfFileExists{alphabeta.sty}{
              \usepackage{alphabeta}
          }{
              \usepackage[mathletters]{ucs}
              \usepackage[utf8x]{inputenc}
          }
    \else
        \usepackage{fontspec}
        \usepackage{unicode-math}
    \fi

    \usepackage{fancyvrb} % verbatim replacement that allows latex
    \usepackage{grffile} % extends the file name processing of package graphics
                         % to support a larger range
    \makeatletter % fix for old versions of grffile with XeLaTeX
    \@ifpackagelater{grffile}{2019/11/01}
    {
      % Do nothing on new versions
    }
    {
      \def\Gread@@xetex#1{%
        \IfFileExists{"\Gin@base".bb}%
        {\Gread@eps{\Gin@base.bb}}%
        {\Gread@@xetex@aux#1}%
      }
    }
    \makeatother
    \usepackage[Export]{adjustbox} % Used to constrain images to a maximum size
    \adjustboxset{max size={0.9\linewidth}{0.9\paperheight}}

    % The hyperref package gives us a pdf with properly built
    % internal navigation ('pdf bookmarks' for the table of contents,
    % internal cross-reference links, web links for URLs, etc.)
    \usepackage{hyperref}
    % The default LaTeX title has an obnoxious amount of whitespace. By default,
    % titling removes some of it. It also provides customization options.
    \usepackage{titling}
    \usepackage{longtable} % longtable support required by pandoc >1.10
    \usepackage{booktabs}  % table support for pandoc > 1.12.2
    \usepackage{array}     % table support for pandoc >= 2.11.3
    \usepackage{calc}      % table minipage width calculation for pandoc >= 2.11.1
    \usepackage[inline]{enumitem} % IRkernel/repr support (it uses the enumerate* environment)
    \usepackage[normalem]{ulem} % ulem is needed to support strikethroughs (\sout)
                                % normalem makes italics be italics, not underlines
    \usepackage{soul}      % strikethrough (\st) support for pandoc >= 3.0.0
    \usepackage{mathrsfs}
    

    
    % Colors for the hyperref package
    \definecolor{urlcolor}{rgb}{0,.145,.698}
    \definecolor{linkcolor}{rgb}{.71,0.21,0.01}
    \definecolor{citecolor}{rgb}{.12,.54,.11}

    % ANSI colors
    \definecolor{ansi-black}{HTML}{3E424D}
    \definecolor{ansi-black-intense}{HTML}{282C36}
    \definecolor{ansi-red}{HTML}{E75C58}
    \definecolor{ansi-red-intense}{HTML}{B22B31}
    \definecolor{ansi-green}{HTML}{00A250}
    \definecolor{ansi-green-intense}{HTML}{007427}
    \definecolor{ansi-yellow}{HTML}{DDB62B}
    \definecolor{ansi-yellow-intense}{HTML}{B27D12}
    \definecolor{ansi-blue}{HTML}{208FFB}
    \definecolor{ansi-blue-intense}{HTML}{0065CA}
    \definecolor{ansi-magenta}{HTML}{D160C4}
    \definecolor{ansi-magenta-intense}{HTML}{A03196}
    \definecolor{ansi-cyan}{HTML}{60C6C8}
    \definecolor{ansi-cyan-intense}{HTML}{258F8F}
    \definecolor{ansi-white}{HTML}{C5C1B4}
    \definecolor{ansi-white-intense}{HTML}{A1A6B2}
    \definecolor{ansi-default-inverse-fg}{HTML}{FFFFFF}
    \definecolor{ansi-default-inverse-bg}{HTML}{000000}

    % common color for the border for error outputs.
    \definecolor{outerrorbackground}{HTML}{FFDFDF}

    % commands and environments needed by pandoc snippets
    % extracted from the output of `pandoc -s`
    \providecommand{\tightlist}{%
      \setlength{\itemsep}{0pt}\setlength{\parskip}{0pt}}
    \DefineVerbatimEnvironment{Highlighting}{Verbatim}{commandchars=\\\{\}}
    % Add ',fontsize=\small' for more characters per line
    \newenvironment{Shaded}{}{}
    \newcommand{\KeywordTok}[1]{\textcolor[rgb]{0.00,0.44,0.13}{\textbf{{#1}}}}
    \newcommand{\DataTypeTok}[1]{\textcolor[rgb]{0.56,0.13,0.00}{{#1}}}
    \newcommand{\DecValTok}[1]{\textcolor[rgb]{0.25,0.63,0.44}{{#1}}}
    \newcommand{\BaseNTok}[1]{\textcolor[rgb]{0.25,0.63,0.44}{{#1}}}
    \newcommand{\FloatTok}[1]{\textcolor[rgb]{0.25,0.63,0.44}{{#1}}}
    \newcommand{\CharTok}[1]{\textcolor[rgb]{0.25,0.44,0.63}{{#1}}}
    \newcommand{\StringTok}[1]{\textcolor[rgb]{0.25,0.44,0.63}{{#1}}}
    \newcommand{\CommentTok}[1]{\textcolor[rgb]{0.38,0.63,0.69}{\textit{{#1}}}}
    \newcommand{\OtherTok}[1]{\textcolor[rgb]{0.00,0.44,0.13}{{#1}}}
    \newcommand{\AlertTok}[1]{\textcolor[rgb]{1.00,0.00,0.00}{\textbf{{#1}}}}
    \newcommand{\FunctionTok}[1]{\textcolor[rgb]{0.02,0.16,0.49}{{#1}}}
    \newcommand{\RegionMarkerTok}[1]{{#1}}
    \newcommand{\ErrorTok}[1]{\textcolor[rgb]{1.00,0.00,0.00}{\textbf{{#1}}}}
    \newcommand{\NormalTok}[1]{{#1}}

    % Additional commands for more recent versions of Pandoc
    \newcommand{\ConstantTok}[1]{\textcolor[rgb]{0.53,0.00,0.00}{{#1}}}
    \newcommand{\SpecialCharTok}[1]{\textcolor[rgb]{0.25,0.44,0.63}{{#1}}}
    \newcommand{\VerbatimStringTok}[1]{\textcolor[rgb]{0.25,0.44,0.63}{{#1}}}
    \newcommand{\SpecialStringTok}[1]{\textcolor[rgb]{0.73,0.40,0.53}{{#1}}}
    \newcommand{\ImportTok}[1]{{#1}}
    \newcommand{\DocumentationTok}[1]{\textcolor[rgb]{0.73,0.13,0.13}{\textit{{#1}}}}
    \newcommand{\AnnotationTok}[1]{\textcolor[rgb]{0.38,0.63,0.69}{\textbf{\textit{{#1}}}}}
    \newcommand{\CommentVarTok}[1]{\textcolor[rgb]{0.38,0.63,0.69}{\textbf{\textit{{#1}}}}}
    \newcommand{\VariableTok}[1]{\textcolor[rgb]{0.10,0.09,0.49}{{#1}}}
    \newcommand{\ControlFlowTok}[1]{\textcolor[rgb]{0.00,0.44,0.13}{\textbf{{#1}}}}
    \newcommand{\OperatorTok}[1]{\textcolor[rgb]{0.40,0.40,0.40}{{#1}}}
    \newcommand{\BuiltInTok}[1]{{#1}}
    \newcommand{\ExtensionTok}[1]{{#1}}
    \newcommand{\PreprocessorTok}[1]{\textcolor[rgb]{0.74,0.48,0.00}{{#1}}}
    \newcommand{\AttributeTok}[1]{\textcolor[rgb]{0.49,0.56,0.16}{{#1}}}
    \newcommand{\InformationTok}[1]{\textcolor[rgb]{0.38,0.63,0.69}{\textbf{\textit{{#1}}}}}
    \newcommand{\WarningTok}[1]{\textcolor[rgb]{0.38,0.63,0.69}{\textbf{\textit{{#1}}}}}
    \makeatletter
    \newsavebox\pandoc@box
    \newcommand*\pandocbounded[1]{%
      \sbox\pandoc@box{#1}%
      % scaling factors for width and height
      \Gscale@div\@tempa\textheight{\dimexpr\ht\pandoc@box+\dp\pandoc@box\relax}%
      \Gscale@div\@tempb\linewidth{\wd\pandoc@box}%
      % select the smaller of both
      \ifdim\@tempb\p@<\@tempa\p@
        \let\@tempa\@tempb
      \fi
      % scaling accordingly (\@tempa < 1)
      \ifdim\@tempa\p@<\p@
        \scalebox{\@tempa}{\usebox\pandoc@box}%
      % scaling not needed, use as it is
      \else
        \usebox{\pandoc@box}%
      \fi
    }
    \makeatother

    % Define a nice break command that doesn't care if a line doesn't already
    % exist.
    \def\br{\hspace*{\fill} \\* }
    % Math Jax compatibility definitions
    \def\gt{>}
    \def\lt{<}
    \let\Oldtex\TeX
    \let\Oldlatex\LaTeX
    \renewcommand{\TeX}{\textrm{\Oldtex}}
    \renewcommand{\LaTeX}{\textrm{\Oldlatex}}
    % Document parameters
    % Document title
    \title{Ch04-classification-lab}
    
    
    
    
    
    
    
% Pygments definitions
\makeatletter
\def\PY@reset{\let\PY@it=\relax \let\PY@bf=\relax%
    \let\PY@ul=\relax \let\PY@tc=\relax%
    \let\PY@bc=\relax \let\PY@ff=\relax}
\def\PY@tok#1{\csname PY@tok@#1\endcsname}
\def\PY@toks#1+{\ifx\relax#1\empty\else%
    \PY@tok{#1}\expandafter\PY@toks\fi}
\def\PY@do#1{\PY@bc{\PY@tc{\PY@ul{%
    \PY@it{\PY@bf{\PY@ff{#1}}}}}}}
\def\PY#1#2{\PY@reset\PY@toks#1+\relax+\PY@do{#2}}

\@namedef{PY@tok@w}{\def\PY@tc##1{\textcolor[rgb]{0.73,0.73,0.73}{##1}}}
\@namedef{PY@tok@c}{\let\PY@it=\textit\def\PY@tc##1{\textcolor[rgb]{0.24,0.48,0.48}{##1}}}
\@namedef{PY@tok@cp}{\def\PY@tc##1{\textcolor[rgb]{0.61,0.40,0.00}{##1}}}
\@namedef{PY@tok@k}{\let\PY@bf=\textbf\def\PY@tc##1{\textcolor[rgb]{0.00,0.50,0.00}{##1}}}
\@namedef{PY@tok@kp}{\def\PY@tc##1{\textcolor[rgb]{0.00,0.50,0.00}{##1}}}
\@namedef{PY@tok@kt}{\def\PY@tc##1{\textcolor[rgb]{0.69,0.00,0.25}{##1}}}
\@namedef{PY@tok@o}{\def\PY@tc##1{\textcolor[rgb]{0.40,0.40,0.40}{##1}}}
\@namedef{PY@tok@ow}{\let\PY@bf=\textbf\def\PY@tc##1{\textcolor[rgb]{0.67,0.13,1.00}{##1}}}
\@namedef{PY@tok@nb}{\def\PY@tc##1{\textcolor[rgb]{0.00,0.50,0.00}{##1}}}
\@namedef{PY@tok@nf}{\def\PY@tc##1{\textcolor[rgb]{0.00,0.00,1.00}{##1}}}
\@namedef{PY@tok@nc}{\let\PY@bf=\textbf\def\PY@tc##1{\textcolor[rgb]{0.00,0.00,1.00}{##1}}}
\@namedef{PY@tok@nn}{\let\PY@bf=\textbf\def\PY@tc##1{\textcolor[rgb]{0.00,0.00,1.00}{##1}}}
\@namedef{PY@tok@ne}{\let\PY@bf=\textbf\def\PY@tc##1{\textcolor[rgb]{0.80,0.25,0.22}{##1}}}
\@namedef{PY@tok@nv}{\def\PY@tc##1{\textcolor[rgb]{0.10,0.09,0.49}{##1}}}
\@namedef{PY@tok@no}{\def\PY@tc##1{\textcolor[rgb]{0.53,0.00,0.00}{##1}}}
\@namedef{PY@tok@nl}{\def\PY@tc##1{\textcolor[rgb]{0.46,0.46,0.00}{##1}}}
\@namedef{PY@tok@ni}{\let\PY@bf=\textbf\def\PY@tc##1{\textcolor[rgb]{0.44,0.44,0.44}{##1}}}
\@namedef{PY@tok@na}{\def\PY@tc##1{\textcolor[rgb]{0.41,0.47,0.13}{##1}}}
\@namedef{PY@tok@nt}{\let\PY@bf=\textbf\def\PY@tc##1{\textcolor[rgb]{0.00,0.50,0.00}{##1}}}
\@namedef{PY@tok@nd}{\def\PY@tc##1{\textcolor[rgb]{0.67,0.13,1.00}{##1}}}
\@namedef{PY@tok@s}{\def\PY@tc##1{\textcolor[rgb]{0.73,0.13,0.13}{##1}}}
\@namedef{PY@tok@sd}{\let\PY@it=\textit\def\PY@tc##1{\textcolor[rgb]{0.73,0.13,0.13}{##1}}}
\@namedef{PY@tok@si}{\let\PY@bf=\textbf\def\PY@tc##1{\textcolor[rgb]{0.64,0.35,0.47}{##1}}}
\@namedef{PY@tok@se}{\let\PY@bf=\textbf\def\PY@tc##1{\textcolor[rgb]{0.67,0.36,0.12}{##1}}}
\@namedef{PY@tok@sr}{\def\PY@tc##1{\textcolor[rgb]{0.64,0.35,0.47}{##1}}}
\@namedef{PY@tok@ss}{\def\PY@tc##1{\textcolor[rgb]{0.10,0.09,0.49}{##1}}}
\@namedef{PY@tok@sx}{\def\PY@tc##1{\textcolor[rgb]{0.00,0.50,0.00}{##1}}}
\@namedef{PY@tok@m}{\def\PY@tc##1{\textcolor[rgb]{0.40,0.40,0.40}{##1}}}
\@namedef{PY@tok@gh}{\let\PY@bf=\textbf\def\PY@tc##1{\textcolor[rgb]{0.00,0.00,0.50}{##1}}}
\@namedef{PY@tok@gu}{\let\PY@bf=\textbf\def\PY@tc##1{\textcolor[rgb]{0.50,0.00,0.50}{##1}}}
\@namedef{PY@tok@gd}{\def\PY@tc##1{\textcolor[rgb]{0.63,0.00,0.00}{##1}}}
\@namedef{PY@tok@gi}{\def\PY@tc##1{\textcolor[rgb]{0.00,0.52,0.00}{##1}}}
\@namedef{PY@tok@gr}{\def\PY@tc##1{\textcolor[rgb]{0.89,0.00,0.00}{##1}}}
\@namedef{PY@tok@ge}{\let\PY@it=\textit}
\@namedef{PY@tok@gs}{\let\PY@bf=\textbf}
\@namedef{PY@tok@ges}{\let\PY@bf=\textbf\let\PY@it=\textit}
\@namedef{PY@tok@gp}{\let\PY@bf=\textbf\def\PY@tc##1{\textcolor[rgb]{0.00,0.00,0.50}{##1}}}
\@namedef{PY@tok@go}{\def\PY@tc##1{\textcolor[rgb]{0.44,0.44,0.44}{##1}}}
\@namedef{PY@tok@gt}{\def\PY@tc##1{\textcolor[rgb]{0.00,0.27,0.87}{##1}}}
\@namedef{PY@tok@err}{\def\PY@bc##1{{\setlength{\fboxsep}{\string -\fboxrule}\fcolorbox[rgb]{1.00,0.00,0.00}{1,1,1}{\strut ##1}}}}
\@namedef{PY@tok@kc}{\let\PY@bf=\textbf\def\PY@tc##1{\textcolor[rgb]{0.00,0.50,0.00}{##1}}}
\@namedef{PY@tok@kd}{\let\PY@bf=\textbf\def\PY@tc##1{\textcolor[rgb]{0.00,0.50,0.00}{##1}}}
\@namedef{PY@tok@kn}{\let\PY@bf=\textbf\def\PY@tc##1{\textcolor[rgb]{0.00,0.50,0.00}{##1}}}
\@namedef{PY@tok@kr}{\let\PY@bf=\textbf\def\PY@tc##1{\textcolor[rgb]{0.00,0.50,0.00}{##1}}}
\@namedef{PY@tok@bp}{\def\PY@tc##1{\textcolor[rgb]{0.00,0.50,0.00}{##1}}}
\@namedef{PY@tok@fm}{\def\PY@tc##1{\textcolor[rgb]{0.00,0.00,1.00}{##1}}}
\@namedef{PY@tok@vc}{\def\PY@tc##1{\textcolor[rgb]{0.10,0.09,0.49}{##1}}}
\@namedef{PY@tok@vg}{\def\PY@tc##1{\textcolor[rgb]{0.10,0.09,0.49}{##1}}}
\@namedef{PY@tok@vi}{\def\PY@tc##1{\textcolor[rgb]{0.10,0.09,0.49}{##1}}}
\@namedef{PY@tok@vm}{\def\PY@tc##1{\textcolor[rgb]{0.10,0.09,0.49}{##1}}}
\@namedef{PY@tok@sa}{\def\PY@tc##1{\textcolor[rgb]{0.73,0.13,0.13}{##1}}}
\@namedef{PY@tok@sb}{\def\PY@tc##1{\textcolor[rgb]{0.73,0.13,0.13}{##1}}}
\@namedef{PY@tok@sc}{\def\PY@tc##1{\textcolor[rgb]{0.73,0.13,0.13}{##1}}}
\@namedef{PY@tok@dl}{\def\PY@tc##1{\textcolor[rgb]{0.73,0.13,0.13}{##1}}}
\@namedef{PY@tok@s2}{\def\PY@tc##1{\textcolor[rgb]{0.73,0.13,0.13}{##1}}}
\@namedef{PY@tok@sh}{\def\PY@tc##1{\textcolor[rgb]{0.73,0.13,0.13}{##1}}}
\@namedef{PY@tok@s1}{\def\PY@tc##1{\textcolor[rgb]{0.73,0.13,0.13}{##1}}}
\@namedef{PY@tok@mb}{\def\PY@tc##1{\textcolor[rgb]{0.40,0.40,0.40}{##1}}}
\@namedef{PY@tok@mf}{\def\PY@tc##1{\textcolor[rgb]{0.40,0.40,0.40}{##1}}}
\@namedef{PY@tok@mh}{\def\PY@tc##1{\textcolor[rgb]{0.40,0.40,0.40}{##1}}}
\@namedef{PY@tok@mi}{\def\PY@tc##1{\textcolor[rgb]{0.40,0.40,0.40}{##1}}}
\@namedef{PY@tok@il}{\def\PY@tc##1{\textcolor[rgb]{0.40,0.40,0.40}{##1}}}
\@namedef{PY@tok@mo}{\def\PY@tc##1{\textcolor[rgb]{0.40,0.40,0.40}{##1}}}
\@namedef{PY@tok@ch}{\let\PY@it=\textit\def\PY@tc##1{\textcolor[rgb]{0.24,0.48,0.48}{##1}}}
\@namedef{PY@tok@cm}{\let\PY@it=\textit\def\PY@tc##1{\textcolor[rgb]{0.24,0.48,0.48}{##1}}}
\@namedef{PY@tok@cpf}{\let\PY@it=\textit\def\PY@tc##1{\textcolor[rgb]{0.24,0.48,0.48}{##1}}}
\@namedef{PY@tok@c1}{\let\PY@it=\textit\def\PY@tc##1{\textcolor[rgb]{0.24,0.48,0.48}{##1}}}
\@namedef{PY@tok@cs}{\let\PY@it=\textit\def\PY@tc##1{\textcolor[rgb]{0.24,0.48,0.48}{##1}}}

\def\PYZbs{\char`\\}
\def\PYZus{\char`\_}
\def\PYZob{\char`\{}
\def\PYZcb{\char`\}}
\def\PYZca{\char`\^}
\def\PYZam{\char`\&}
\def\PYZlt{\char`\<}
\def\PYZgt{\char`\>}
\def\PYZsh{\char`\#}
\def\PYZpc{\char`\%}
\def\PYZdl{\char`\$}
\def\PYZhy{\char`\-}
\def\PYZsq{\char`\'}
\def\PYZdq{\char`\"}
\def\PYZti{\char`\~}
% for compatibility with earlier versions
\def\PYZat{@}
\def\PYZlb{[}
\def\PYZrb{]}
\makeatother


    % For linebreaks inside Verbatim environment from package fancyvrb.
    \makeatletter
        \newbox\Wrappedcontinuationbox
        \newbox\Wrappedvisiblespacebox
        \newcommand*\Wrappedvisiblespace {\textcolor{red}{\textvisiblespace}}
        \newcommand*\Wrappedcontinuationsymbol {\textcolor{red}{\llap{\tiny$\m@th\hookrightarrow$}}}
        \newcommand*\Wrappedcontinuationindent {3ex }
        \newcommand*\Wrappedafterbreak {\kern\Wrappedcontinuationindent\copy\Wrappedcontinuationbox}
        % Take advantage of the already applied Pygments mark-up to insert
        % potential linebreaks for TeX processing.
        %        {, <, #, %, $, ' and ": go to next line.
        %        _, }, ^, &, >, - and ~: stay at end of broken line.
        % Use of \textquotesingle for straight quote.
        \newcommand*\Wrappedbreaksatspecials {%
            \def\PYGZus{\discretionary{\char`\_}{\Wrappedafterbreak}{\char`\_}}%
            \def\PYGZob{\discretionary{}{\Wrappedafterbreak\char`\{}{\char`\{}}%
            \def\PYGZcb{\discretionary{\char`\}}{\Wrappedafterbreak}{\char`\}}}%
            \def\PYGZca{\discretionary{\char`\^}{\Wrappedafterbreak}{\char`\^}}%
            \def\PYGZam{\discretionary{\char`\&}{\Wrappedafterbreak}{\char`\&}}%
            \def\PYGZlt{\discretionary{}{\Wrappedafterbreak\char`\<}{\char`\<}}%
            \def\PYGZgt{\discretionary{\char`\>}{\Wrappedafterbreak}{\char`\>}}%
            \def\PYGZsh{\discretionary{}{\Wrappedafterbreak\char`\#}{\char`\#}}%
            \def\PYGZpc{\discretionary{}{\Wrappedafterbreak\char`\%}{\char`\%}}%
            \def\PYGZdl{\discretionary{}{\Wrappedafterbreak\char`\$}{\char`\$}}%
            \def\PYGZhy{\discretionary{\char`\-}{\Wrappedafterbreak}{\char`\-}}%
            \def\PYGZsq{\discretionary{}{\Wrappedafterbreak\textquotesingle}{\textquotesingle}}%
            \def\PYGZdq{\discretionary{}{\Wrappedafterbreak\char`\"}{\char`\"}}%
            \def\PYGZti{\discretionary{\char`\~}{\Wrappedafterbreak}{\char`\~}}%
        }
        % Some characters . , ; ? ! / are not pygmentized.
        % This macro makes them "active" and they will insert potential linebreaks
        \newcommand*\Wrappedbreaksatpunct {%
            \lccode`\~`\.\lowercase{\def~}{\discretionary{\hbox{\char`\.}}{\Wrappedafterbreak}{\hbox{\char`\.}}}%
            \lccode`\~`\,\lowercase{\def~}{\discretionary{\hbox{\char`\,}}{\Wrappedafterbreak}{\hbox{\char`\,}}}%
            \lccode`\~`\;\lowercase{\def~}{\discretionary{\hbox{\char`\;}}{\Wrappedafterbreak}{\hbox{\char`\;}}}%
            \lccode`\~`\:\lowercase{\def~}{\discretionary{\hbox{\char`\:}}{\Wrappedafterbreak}{\hbox{\char`\:}}}%
            \lccode`\~`\?\lowercase{\def~}{\discretionary{\hbox{\char`\?}}{\Wrappedafterbreak}{\hbox{\char`\?}}}%
            \lccode`\~`\!\lowercase{\def~}{\discretionary{\hbox{\char`\!}}{\Wrappedafterbreak}{\hbox{\char`\!}}}%
            \lccode`\~`\/\lowercase{\def~}{\discretionary{\hbox{\char`\/}}{\Wrappedafterbreak}{\hbox{\char`\/}}}%
            \catcode`\.\active
            \catcode`\,\active
            \catcode`\;\active
            \catcode`\:\active
            \catcode`\?\active
            \catcode`\!\active
            \catcode`\/\active
            \lccode`\~`\~
        }
    \makeatother

    \let\OriginalVerbatim=\Verbatim
    \makeatletter
    \renewcommand{\Verbatim}[1][1]{%
        %\parskip\z@skip
        \sbox\Wrappedcontinuationbox {\Wrappedcontinuationsymbol}%
        \sbox\Wrappedvisiblespacebox {\FV@SetupFont\Wrappedvisiblespace}%
        \def\FancyVerbFormatLine ##1{\hsize\linewidth
            \vtop{\raggedright\hyphenpenalty\z@\exhyphenpenalty\z@
                \doublehyphendemerits\z@\finalhyphendemerits\z@
                \strut ##1\strut}%
        }%
        % If the linebreak is at a space, the latter will be displayed as visible
        % space at end of first line, and a continuation symbol starts next line.
        % Stretch/shrink are however usually zero for typewriter font.
        \def\FV@Space {%
            \nobreak\hskip\z@ plus\fontdimen3\font minus\fontdimen4\font
            \discretionary{\copy\Wrappedvisiblespacebox}{\Wrappedafterbreak}
            {\kern\fontdimen2\font}%
        }%

        % Allow breaks at special characters using \PYG... macros.
        \Wrappedbreaksatspecials
        % Breaks at punctuation characters . , ; ? ! and / need catcode=\active
        \OriginalVerbatim[#1,codes*=\Wrappedbreaksatpunct]%
    }
    \makeatother

    % Exact colors from NB
    \definecolor{incolor}{HTML}{303F9F}
    \definecolor{outcolor}{HTML}{D84315}
    \definecolor{cellborder}{HTML}{CFCFCF}
    \definecolor{cellbackground}{HTML}{F7F7F7}

    % prompt
    \makeatletter
    \newcommand{\boxspacing}{\kern\kvtcb@left@rule\kern\kvtcb@boxsep}
    \makeatother
    \newcommand{\prompt}[4]{
        {\ttfamily\llap{{\color{#2}[#3]:\hspace{3pt}#4}}\vspace{-\baselineskip}}
    }
    

    
    % Prevent overflowing lines due to hard-to-break entities
    \sloppy
    % Setup hyperref package
    \hypersetup{
      breaklinks=true,  % so long urls are correctly broken across lines
      colorlinks=true,
      urlcolor=urlcolor,
      linkcolor=linkcolor,
      citecolor=citecolor,
      }
    % Slightly bigger margins than the latex defaults
    
    \geometry{verbose,tmargin=1in,bmargin=1in,lmargin=1in,rmargin=1in}
    
    

\begin{document}
    
    \maketitle
    
    

    
    \hypertarget{logistic-regression-lda-qda-and-knn}{%
\section{Logistic Regression, LDA, QDA, and
KNN}\label{logistic-regression-lda-qda-and-knn}}

    \hypertarget{the-stock-market-data}{%
\subsection{The Stock Market Data}\label{the-stock-market-data}}

In this lab we will examine the \texttt{Smarket} data, which is part of
the \texttt{ISLP} library. This data set consists of percentage returns
for the S\&P 500 stock index over 1,250 days, from the beginning of 2001
until the end of 2005. For each date, we have recorded the percentage
returns for each of the five previous trading days, \texttt{Lag1}
through \texttt{Lag5}. We have also recorded \texttt{Volume} (the number
of shares traded on the previous day, in billions), \texttt{Today} (the
percentage return on the date in question) and \texttt{Direction}
(whether the market was \texttt{Up} or \texttt{Down} on this date).

We start by importing our libraries at this top level; these are all
imports we have seen in previous labs.

    \begin{tcolorbox}[breakable, size=fbox, boxrule=1pt, pad at break*=1mm,colback=cellbackground, colframe=cellborder]
\prompt{In}{incolor}{1}{\boxspacing}
\begin{Verbatim}[commandchars=\\\{\}]
\PY{k+kn}{import}\PY{+w}{ }\PY{n+nn}{numpy}\PY{+w}{ }\PY{k}{as}\PY{+w}{ }\PY{n+nn}{np}
\PY{k+kn}{import}\PY{+w}{ }\PY{n+nn}{pandas}\PY{+w}{ }\PY{k}{as}\PY{+w}{ }\PY{n+nn}{pd}
\PY{k+kn}{from}\PY{+w}{ }\PY{n+nn}{matplotlib}\PY{n+nn}{.}\PY{n+nn}{pyplot}\PY{+w}{ }\PY{k+kn}{import} \PY{n}{subplots}
\PY{k+kn}{import}\PY{+w}{ }\PY{n+nn}{statsmodels}\PY{n+nn}{.}\PY{n+nn}{api}\PY{+w}{ }\PY{k}{as}\PY{+w}{ }\PY{n+nn}{sm}
\PY{k+kn}{from}\PY{+w}{ }\PY{n+nn}{ISLP}\PY{+w}{ }\PY{k+kn}{import} \PY{n}{load\PYZus{}data}
\PY{k+kn}{from}\PY{+w}{ }\PY{n+nn}{ISLP}\PY{n+nn}{.}\PY{n+nn}{models}\PY{+w}{ }\PY{k+kn}{import} \PY{p}{(}\PY{n}{ModelSpec} \PY{k}{as} \PY{n}{MS}\PY{p}{,}
                         \PY{n}{summarize}\PY{p}{)}
\end{Verbatim}
\end{tcolorbox}

    We also collect together the new imports needed for this lab.

    \begin{tcolorbox}[breakable, size=fbox, boxrule=1pt, pad at break*=1mm,colback=cellbackground, colframe=cellborder]
\prompt{In}{incolor}{2}{\boxspacing}
\begin{Verbatim}[commandchars=\\\{\}]
\PY{k+kn}{from}\PY{+w}{ }\PY{n+nn}{ISLP}\PY{+w}{ }\PY{k+kn}{import} \PY{n}{confusion\PYZus{}table}
\PY{k+kn}{from}\PY{+w}{ }\PY{n+nn}{ISLP}\PY{n+nn}{.}\PY{n+nn}{models}\PY{+w}{ }\PY{k+kn}{import} \PY{n}{contrast}
\PY{k+kn}{from}\PY{+w}{ }\PY{n+nn}{sklearn}\PY{n+nn}{.}\PY{n+nn}{discriminant\PYZus{}analysis}\PY{+w}{ }\PY{k+kn}{import} \PYZbs{}
     \PY{p}{(}\PY{n}{LinearDiscriminantAnalysis} \PY{k}{as} \PY{n}{LDA}\PY{p}{,}
      \PY{n}{QuadraticDiscriminantAnalysis} \PY{k}{as} \PY{n}{QDA}\PY{p}{)}
\PY{k+kn}{from}\PY{+w}{ }\PY{n+nn}{sklearn}\PY{n+nn}{.}\PY{n+nn}{naive\PYZus{}bayes}\PY{+w}{ }\PY{k+kn}{import} \PY{n}{GaussianNB}
\PY{k+kn}{from}\PY{+w}{ }\PY{n+nn}{sklearn}\PY{n+nn}{.}\PY{n+nn}{neighbors}\PY{+w}{ }\PY{k+kn}{import} \PY{n}{KNeighborsClassifier}
\PY{k+kn}{from}\PY{+w}{ }\PY{n+nn}{sklearn}\PY{n+nn}{.}\PY{n+nn}{preprocessing}\PY{+w}{ }\PY{k+kn}{import} \PY{n}{StandardScaler}
\PY{k+kn}{from}\PY{+w}{ }\PY{n+nn}{sklearn}\PY{n+nn}{.}\PY{n+nn}{model\PYZus{}selection}\PY{+w}{ }\PY{k+kn}{import} \PY{n}{train\PYZus{}test\PYZus{}split}
\PY{k+kn}{from}\PY{+w}{ }\PY{n+nn}{sklearn}\PY{n+nn}{.}\PY{n+nn}{linear\PYZus{}model}\PY{+w}{ }\PY{k+kn}{import} \PY{n}{LogisticRegression}
\end{Verbatim}
\end{tcolorbox}

    Now we are ready to load the \texttt{Smarket} data.

    \begin{tcolorbox}[breakable, size=fbox, boxrule=1pt, pad at break*=1mm,colback=cellbackground, colframe=cellborder]
\prompt{In}{incolor}{3}{\boxspacing}
\begin{Verbatim}[commandchars=\\\{\}]
\PY{n}{Smarket} \PY{o}{=} \PY{n}{load\PYZus{}data}\PY{p}{(}\PY{l+s+s1}{\PYZsq{}}\PY{l+s+s1}{Smarket}\PY{l+s+s1}{\PYZsq{}}\PY{p}{)}
\PY{n}{Smarket}
\end{Verbatim}
\end{tcolorbox}

            \begin{tcolorbox}[breakable, size=fbox, boxrule=.5pt, pad at break*=1mm, opacityfill=0]
\prompt{Out}{outcolor}{3}{\boxspacing}
\begin{Verbatim}[commandchars=\\\{\}]
      Year   Lag1   Lag2   Lag3   Lag4   Lag5   Volume  Today Direction
0     2001  0.381 -0.192 -2.624 -1.055  5.010  1.19130  0.959        Up
1     2001  0.959  0.381 -0.192 -2.624 -1.055  1.29650  1.032        Up
2     2001  1.032  0.959  0.381 -0.192 -2.624  1.41120 -0.623      Down
3     2001 -0.623  1.032  0.959  0.381 -0.192  1.27600  0.614        Up
4     2001  0.614 -0.623  1.032  0.959  0.381  1.20570  0.213        Up
{\ldots}    {\ldots}    {\ldots}    {\ldots}    {\ldots}    {\ldots}    {\ldots}      {\ldots}    {\ldots}       {\ldots}
1245  2005  0.422  0.252 -0.024 -0.584 -0.285  1.88850  0.043        Up
1246  2005  0.043  0.422  0.252 -0.024 -0.584  1.28581 -0.955      Down
1247  2005 -0.955  0.043  0.422  0.252 -0.024  1.54047  0.130        Up
1248  2005  0.130 -0.955  0.043  0.422  0.252  1.42236 -0.298      Down
1249  2005 -0.298  0.130 -0.955  0.043  0.422  1.38254 -0.489      Down

[1250 rows x 9 columns]
\end{Verbatim}
\end{tcolorbox}
        
    This gives a truncated listing of the data. We can see what the variable
names are.

    \begin{tcolorbox}[breakable, size=fbox, boxrule=1pt, pad at break*=1mm,colback=cellbackground, colframe=cellborder]
\prompt{In}{incolor}{4}{\boxspacing}
\begin{Verbatim}[commandchars=\\\{\}]
\PY{n}{Smarket}\PY{o}{.}\PY{n}{columns}
\end{Verbatim}
\end{tcolorbox}

            \begin{tcolorbox}[breakable, size=fbox, boxrule=.5pt, pad at break*=1mm, opacityfill=0]
\prompt{Out}{outcolor}{4}{\boxspacing}
\begin{Verbatim}[commandchars=\\\{\}]
Index(['Year', 'Lag1', 'Lag2', 'Lag3', 'Lag4', 'Lag5', 'Volume', 'Today',
       'Direction'],
      dtype='object')
\end{Verbatim}
\end{tcolorbox}
        
    We compute the correlation matrix using the \texttt{corr()} method for
data frames, which produces a matrix that contains all of the pairwise
correlations among the variables.

By instructing \texttt{pandas} to use only numeric variables, the
\texttt{corr()} method does not report a correlation for the
\texttt{Direction} variable because it is qualitative.

    \begin{tcolorbox}[breakable, size=fbox, boxrule=1pt, pad at break*=1mm,colback=cellbackground, colframe=cellborder]
\prompt{In}{incolor}{5}{\boxspacing}
\begin{Verbatim}[commandchars=\\\{\}]
\PY{n}{Smarket}\PY{o}{.}\PY{n}{corr}\PY{p}{(}\PY{n}{numeric\PYZus{}only}\PY{o}{=}\PY{k+kc}{True}\PY{p}{)}
\end{Verbatim}
\end{tcolorbox}

            \begin{tcolorbox}[breakable, size=fbox, boxrule=.5pt, pad at break*=1mm, opacityfill=0]
\prompt{Out}{outcolor}{5}{\boxspacing}
\begin{Verbatim}[commandchars=\\\{\}]
            Year      Lag1      Lag2      Lag3      Lag4      Lag5    Volume  \textbackslash{}
Year    1.000000  0.029700  0.030596  0.033195  0.035689  0.029788  0.539006
Lag1    0.029700  1.000000 -0.026294 -0.010803 -0.002986 -0.005675  0.040910
Lag2    0.030596 -0.026294  1.000000 -0.025897 -0.010854 -0.003558 -0.043383
Lag3    0.033195 -0.010803 -0.025897  1.000000 -0.024051 -0.018808 -0.041824
Lag4    0.035689 -0.002986 -0.010854 -0.024051  1.000000 -0.027084 -0.048414
Lag5    0.029788 -0.005675 -0.003558 -0.018808 -0.027084  1.000000 -0.022002
Volume  0.539006  0.040910 -0.043383 -0.041824 -0.048414 -0.022002  1.000000
Today   0.030095 -0.026155 -0.010250 -0.002448 -0.006900 -0.034860  0.014592

           Today
Year    0.030095
Lag1   -0.026155
Lag2   -0.010250
Lag3   -0.002448
Lag4   -0.006900
Lag5   -0.034860
Volume  0.014592
Today   1.000000
\end{Verbatim}
\end{tcolorbox}
        
    As one would expect, the correlations between the lagged return
variables and today's return are close to zero. The only substantial
correlation is between \texttt{Year} and \texttt{Volume}. By plotting
the data we see that \texttt{Volume} is increasing over time. In other
words, the average number of shares traded daily increased from 2001 to
2005.

    \begin{tcolorbox}[breakable, size=fbox, boxrule=1pt, pad at break*=1mm,colback=cellbackground, colframe=cellborder]
\prompt{In}{incolor}{6}{\boxspacing}
\begin{Verbatim}[commandchars=\\\{\}]
\PY{n}{Smarket}\PY{o}{.}\PY{n}{plot}\PY{p}{(}\PY{n}{y}\PY{o}{=}\PY{l+s+s1}{\PYZsq{}}\PY{l+s+s1}{Volume}\PY{l+s+s1}{\PYZsq{}}\PY{p}{)}\PY{p}{;}
\end{Verbatim}
\end{tcolorbox}

    \begin{center}
    \adjustimage{max size={0.9\linewidth}{0.9\paperheight}}{artifacts/latex/Ch04-classification-lab_files/artifacts/latex/Ch04-classification-lab_12_0.png}
    \end{center}
    { \hspace*{\fill} \\}
    
    \hypertarget{logistic-regression}{%
\subsection{Logistic Regression}\label{logistic-regression}}

Next, we will fit a logistic regression model in order to predict
\texttt{Direction} using \texttt{Lag1} through \texttt{Lag5} and
\texttt{Volume}. The \texttt{sm.GLM()} function fits \emph{generalized
linear models}, a class of models that includes logistic regression.
Alternatively, the function \texttt{sm.Logit()} fits a logistic
regression model directly. The syntax of \texttt{sm.GLM()} is similar to
that of \texttt{sm.OLS()}, except that we must pass in the argument
\texttt{family=sm.families.Binomial()} in order to tell
\texttt{statsmodels} to run a logistic regression rather than some other
type of generalized linear model.

    \begin{tcolorbox}[breakable, size=fbox, boxrule=1pt, pad at break*=1mm,colback=cellbackground, colframe=cellborder]
\prompt{In}{incolor}{7}{\boxspacing}
\begin{Verbatim}[commandchars=\\\{\}]
\PY{n}{allvars} \PY{o}{=} \PY{n}{Smarket}\PY{o}{.}\PY{n}{columns}\PY{o}{.}\PY{n}{drop}\PY{p}{(}\PY{p}{[}\PY{l+s+s1}{\PYZsq{}}\PY{l+s+s1}{Today}\PY{l+s+s1}{\PYZsq{}}\PY{p}{,} \PY{l+s+s1}{\PYZsq{}}\PY{l+s+s1}{Direction}\PY{l+s+s1}{\PYZsq{}}\PY{p}{,} \PY{l+s+s1}{\PYZsq{}}\PY{l+s+s1}{Year}\PY{l+s+s1}{\PYZsq{}}\PY{p}{]}\PY{p}{)}
\PY{n}{design} \PY{o}{=} \PY{n}{MS}\PY{p}{(}\PY{n}{allvars}\PY{p}{)}
\PY{n}{X} \PY{o}{=} \PY{n}{design}\PY{o}{.}\PY{n}{fit\PYZus{}transform}\PY{p}{(}\PY{n}{Smarket}\PY{p}{)}
\PY{n}{y} \PY{o}{=} \PY{n}{Smarket}\PY{o}{.}\PY{n}{Direction} \PY{o}{==} \PY{l+s+s1}{\PYZsq{}}\PY{l+s+s1}{Up}\PY{l+s+s1}{\PYZsq{}}
\PY{n}{glm} \PY{o}{=} \PY{n}{sm}\PY{o}{.}\PY{n}{GLM}\PY{p}{(}\PY{n}{y}\PY{p}{,}
             \PY{n}{X}\PY{p}{,}
             \PY{n}{family}\PY{o}{=}\PY{n}{sm}\PY{o}{.}\PY{n}{families}\PY{o}{.}\PY{n}{Binomial}\PY{p}{(}\PY{p}{)}\PY{p}{)}
\PY{n}{results} \PY{o}{=} \PY{n}{glm}\PY{o}{.}\PY{n}{fit}\PY{p}{(}\PY{p}{)}
\PY{n}{summarize}\PY{p}{(}\PY{n}{results}\PY{p}{)}
\end{Verbatim}
\end{tcolorbox}

            \begin{tcolorbox}[breakable, size=fbox, boxrule=.5pt, pad at break*=1mm, opacityfill=0]
\prompt{Out}{outcolor}{7}{\boxspacing}
\begin{Verbatim}[commandchars=\\\{\}]
             coef  std err      z  P>|z|
intercept -0.1260    0.241 -0.523  0.601
Lag1      -0.0731    0.050 -1.457  0.145
Lag2      -0.0423    0.050 -0.845  0.398
Lag3       0.0111    0.050  0.222  0.824
Lag4       0.0094    0.050  0.187  0.851
Lag5       0.0103    0.050  0.208  0.835
Volume     0.1354    0.158  0.855  0.392
\end{Verbatim}
\end{tcolorbox}
        
    The smallest \emph{p}-value here is associated with \texttt{Lag1}. The
negative coefficient for this predictor suggests that if the market had
a positive return yesterday, then it is less likely to go up today.
However, at a value of 0.15, the \emph{p}-value is still relatively
large, and so there is no clear evidence of a real association between
\texttt{Lag1} and \texttt{Direction}.

We use the \texttt{params} attribute of \texttt{results} in order to
access just the coefficients for this fitted model.

    \begin{tcolorbox}[breakable, size=fbox, boxrule=1pt, pad at break*=1mm,colback=cellbackground, colframe=cellborder]
\prompt{In}{incolor}{8}{\boxspacing}
\begin{Verbatim}[commandchars=\\\{\}]
\PY{n}{results}\PY{o}{.}\PY{n}{params}
\end{Verbatim}
\end{tcolorbox}

            \begin{tcolorbox}[breakable, size=fbox, boxrule=.5pt, pad at break*=1mm, opacityfill=0]
\prompt{Out}{outcolor}{8}{\boxspacing}
\begin{Verbatim}[commandchars=\\\{\}]
intercept   -0.126000
Lag1        -0.073074
Lag2        -0.042301
Lag3         0.011085
Lag4         0.009359
Lag5         0.010313
Volume       0.135441
dtype: float64
\end{Verbatim}
\end{tcolorbox}
        
    Likewise we can use the \texttt{pvalues} attribute to access the
\emph{p}-values for the coefficients.

    \begin{tcolorbox}[breakable, size=fbox, boxrule=1pt, pad at break*=1mm,colback=cellbackground, colframe=cellborder]
\prompt{In}{incolor}{9}{\boxspacing}
\begin{Verbatim}[commandchars=\\\{\}]
\PY{n}{results}\PY{o}{.}\PY{n}{pvalues}
\end{Verbatim}
\end{tcolorbox}

            \begin{tcolorbox}[breakable, size=fbox, boxrule=.5pt, pad at break*=1mm, opacityfill=0]
\prompt{Out}{outcolor}{9}{\boxspacing}
\begin{Verbatim}[commandchars=\\\{\}]
intercept    0.600700
Lag1         0.145232
Lag2         0.398352
Lag3         0.824334
Lag4         0.851445
Lag5         0.834998
Volume       0.392404
dtype: float64
\end{Verbatim}
\end{tcolorbox}
        
    The \texttt{predict()} method of \texttt{results} can be used to predict
the probability that the market will go up, given values of the
predictors. This method returns predictions on the probability scale. If
no data set is supplied to the \texttt{predict()} function, then the
probabilities are computed for the training data that was used to fit
the logistic regression model. As with linear regression, one can pass
an optional \texttt{exog} argument consistent with a design matrix if
desired. Here we have printed only the first ten probabilities.

    \begin{tcolorbox}[breakable, size=fbox, boxrule=1pt, pad at break*=1mm,colback=cellbackground, colframe=cellborder]
\prompt{In}{incolor}{10}{\boxspacing}
\begin{Verbatim}[commandchars=\\\{\}]
\PY{n}{probs} \PY{o}{=} \PY{n}{results}\PY{o}{.}\PY{n}{predict}\PY{p}{(}\PY{p}{)}
\PY{n}{probs}\PY{p}{[}\PY{p}{:}\PY{l+m+mi}{10}\PY{p}{]}
\end{Verbatim}
\end{tcolorbox}

            \begin{tcolorbox}[breakable, size=fbox, boxrule=.5pt, pad at break*=1mm, opacityfill=0]
\prompt{Out}{outcolor}{10}{\boxspacing}
\begin{Verbatim}[commandchars=\\\{\}]
array([0.50708413, 0.48146788, 0.48113883, 0.51522236, 0.51078116,
       0.50695646, 0.49265087, 0.50922916, 0.51761353, 0.48883778])
\end{Verbatim}
\end{tcolorbox}
        
    In order to make a prediction as to whether the market will go up or
down on a particular day, we must convert these predicted probabilities
into class labels, \texttt{Up} or \texttt{Down}. The following two
commands create a vector of class predictions based on whether the
predicted probability of a market increase is greater than or less than
0.5.

    \begin{tcolorbox}[breakable, size=fbox, boxrule=1pt, pad at break*=1mm,colback=cellbackground, colframe=cellborder]
\prompt{In}{incolor}{11}{\boxspacing}
\begin{Verbatim}[commandchars=\\\{\}]
\PY{n}{labels} \PY{o}{=} \PY{n}{np}\PY{o}{.}\PY{n}{array}\PY{p}{(}\PY{p}{[}\PY{l+s+s1}{\PYZsq{}}\PY{l+s+s1}{Down}\PY{l+s+s1}{\PYZsq{}}\PY{p}{]}\PY{o}{*}\PY{l+m+mi}{1250}\PY{p}{)}
\PY{n}{labels}\PY{p}{[}\PY{n}{probs}\PY{o}{\PYZgt{}}\PY{l+m+mf}{0.5}\PY{p}{]} \PY{o}{=} \PY{l+s+s2}{\PYZdq{}}\PY{l+s+s2}{Up}\PY{l+s+s2}{\PYZdq{}}
\end{Verbatim}
\end{tcolorbox}

    The \texttt{confusion\_table()} function from the \texttt{ISLP} package
summarizes these predictions, showing how many observations were
correctly or incorrectly classified. Our function, which is adapted from
a similar function in the module \texttt{sklearn.metrics}, transposes
the resulting matrix and includes row and column labels. The
\texttt{confusion\_table()} function takes as first argument the
predicted labels, and second argument the true labels.

    \begin{tcolorbox}[breakable, size=fbox, boxrule=1pt, pad at break*=1mm,colback=cellbackground, colframe=cellborder]
\prompt{In}{incolor}{12}{\boxspacing}
\begin{Verbatim}[commandchars=\\\{\}]
\PY{n}{confusion\PYZus{}table}\PY{p}{(}\PY{n}{labels}\PY{p}{,} \PY{n}{Smarket}\PY{o}{.}\PY{n}{Direction}\PY{p}{)}
\end{Verbatim}
\end{tcolorbox}

            \begin{tcolorbox}[breakable, size=fbox, boxrule=.5pt, pad at break*=1mm, opacityfill=0]
\prompt{Out}{outcolor}{12}{\boxspacing}
\begin{Verbatim}[commandchars=\\\{\}]
Truth      Down   Up
Predicted
Down        145  141
Up          457  507
\end{Verbatim}
\end{tcolorbox}
        
    The diagonal elements of the confusion matrix indicate correct
predictions, while the off-diagonals represent incorrect predictions.
Hence our model correctly predicted that the market would go up on 507
days and that it would go down on 145 days, for a total of 507 + 145 =
652 correct predictions. The \texttt{np.mean()} function can be used to
compute the fraction of days for which the prediction was correct. In
this case, logistic regression correctly predicted the movement of the
market 52.2\% of the time.

    \begin{tcolorbox}[breakable, size=fbox, boxrule=1pt, pad at break*=1mm,colback=cellbackground, colframe=cellborder]
\prompt{In}{incolor}{13}{\boxspacing}
\begin{Verbatim}[commandchars=\\\{\}]
\PY{p}{(}\PY{l+m+mi}{507}\PY{o}{+}\PY{l+m+mi}{145}\PY{p}{)}\PY{o}{/}\PY{l+m+mi}{1250}\PY{p}{,} \PY{n}{np}\PY{o}{.}\PY{n}{mean}\PY{p}{(}\PY{n}{labels} \PY{o}{==} \PY{n}{Smarket}\PY{o}{.}\PY{n}{Direction}\PY{p}{)}
\end{Verbatim}
\end{tcolorbox}

            \begin{tcolorbox}[breakable, size=fbox, boxrule=.5pt, pad at break*=1mm, opacityfill=0]
\prompt{Out}{outcolor}{13}{\boxspacing}
\begin{Verbatim}[commandchars=\\\{\}]
(0.5216, 0.5216)
\end{Verbatim}
\end{tcolorbox}
        
    

    At first glance, it appears that the logistic regression model is
working a little better than random guessing. However, this result is
misleading because we trained and tested the model on the same set of
1,250 observations. In other words, \(100-52.2=47.8%
\) is the \emph{training} error rate. As we have seen previously, the
training error rate is often overly optimistic --- it tends to
underestimate the test error rate. In order to better assess the
accuracy of the logistic regression model in this setting, we can fit
the model using part of the data, and then examine how well it predicts
the \emph{held out} data. This will yield a more realistic error rate,
in the sense that in practice we will be interested in our model's
performance not on the data that we used to fit the model, but rather on
days in the future for which the market's movements are unknown.

To implement this strategy, we first create a Boolean vector
corresponding to the observations from 2001 through 2004. We then use
this vector to create a held out data set of observations from 2005.

    \begin{tcolorbox}[breakable, size=fbox, boxrule=1pt, pad at break*=1mm,colback=cellbackground, colframe=cellborder]
\prompt{In}{incolor}{14}{\boxspacing}
\begin{Verbatim}[commandchars=\\\{\}]
\PY{n}{train} \PY{o}{=} \PY{p}{(}\PY{n}{Smarket}\PY{o}{.}\PY{n}{Year} \PY{o}{\PYZlt{}} \PY{l+m+mi}{2005}\PY{p}{)}
\PY{n}{Smarket\PYZus{}train} \PY{o}{=} \PY{n}{Smarket}\PY{o}{.}\PY{n}{loc}\PY{p}{[}\PY{n}{train}\PY{p}{]}
\PY{n}{Smarket\PYZus{}test} \PY{o}{=} \PY{n}{Smarket}\PY{o}{.}\PY{n}{loc}\PY{p}{[}\PY{o}{\PYZti{}}\PY{n}{train}\PY{p}{]}
\PY{n}{Smarket\PYZus{}test}\PY{o}{.}\PY{n}{shape}
\end{Verbatim}
\end{tcolorbox}

            \begin{tcolorbox}[breakable, size=fbox, boxrule=.5pt, pad at break*=1mm, opacityfill=0]
\prompt{Out}{outcolor}{14}{\boxspacing}
\begin{Verbatim}[commandchars=\\\{\}]
(252, 9)
\end{Verbatim}
\end{tcolorbox}
        
    The object \texttt{train} is a vector of 1,250 elements, corresponding
to the observations in our data set. The elements of the vector that
correspond to observations that occurred before 2005 are set to
\texttt{True}, whereas those that correspond to observations in 2005 are
set to \texttt{False}. Hence \texttt{train} is a \emph{boolean} array,
since its elements are \texttt{True} and \texttt{False}. Boolean arrays
can be used to obtain a subset of the rows or columns of a data frame
using the \texttt{loc} method. For instance, the command
\texttt{Smarket.loc{[}train{]}} would pick out a submatrix of the stock
market data set, corresponding only to the dates before 2005, since
those are the ones for which the elements of \texttt{train} are
\texttt{True}. The \texttt{\textasciitilde{}} symbol can be used to
negate all of the elements of a Boolean vector. That is,
\texttt{\textasciitilde{}train} is a vector similar to \texttt{train},
except that the elements that are \texttt{True} in \texttt{train} get
swapped to \texttt{False} in \texttt{\textasciitilde{}train}, and vice
versa. Therefore, \texttt{Smarket.loc{[}\textasciitilde{}train{]}}
yields a subset of the rows of the data frame of the stock market data
containing only the observations for which \texttt{train} is
\texttt{False}. The output above indicates that there are 252 such
observations.

We now fit a logistic regression model using only the subset of the
observations that correspond to dates before 2005. We then obtain
predicted probabilities of the stock market going up for each of the
days in our test set --- that is, for the days in 2005.

    \begin{tcolorbox}[breakable, size=fbox, boxrule=1pt, pad at break*=1mm,colback=cellbackground, colframe=cellborder]
\prompt{In}{incolor}{15}{\boxspacing}
\begin{Verbatim}[commandchars=\\\{\}]
\PY{n}{X\PYZus{}train}\PY{p}{,} \PY{n}{X\PYZus{}test} \PY{o}{=} \PY{n}{X}\PY{o}{.}\PY{n}{loc}\PY{p}{[}\PY{n}{train}\PY{p}{]}\PY{p}{,} \PY{n}{X}\PY{o}{.}\PY{n}{loc}\PY{p}{[}\PY{o}{\PYZti{}}\PY{n}{train}\PY{p}{]}
\PY{n}{y\PYZus{}train}\PY{p}{,} \PY{n}{y\PYZus{}test} \PY{o}{=} \PY{n}{y}\PY{o}{.}\PY{n}{loc}\PY{p}{[}\PY{n}{train}\PY{p}{]}\PY{p}{,} \PY{n}{y}\PY{o}{.}\PY{n}{loc}\PY{p}{[}\PY{o}{\PYZti{}}\PY{n}{train}\PY{p}{]}
\PY{n}{glm\PYZus{}train} \PY{o}{=} \PY{n}{sm}\PY{o}{.}\PY{n}{GLM}\PY{p}{(}\PY{n}{y\PYZus{}train}\PY{p}{,}
                   \PY{n}{X\PYZus{}train}\PY{p}{,}
                   \PY{n}{family}\PY{o}{=}\PY{n}{sm}\PY{o}{.}\PY{n}{families}\PY{o}{.}\PY{n}{Binomial}\PY{p}{(}\PY{p}{)}\PY{p}{)}
\PY{n}{results} \PY{o}{=} \PY{n}{glm\PYZus{}train}\PY{o}{.}\PY{n}{fit}\PY{p}{(}\PY{p}{)}
\PY{n}{probs} \PY{o}{=} \PY{n}{results}\PY{o}{.}\PY{n}{predict}\PY{p}{(}\PY{n}{exog}\PY{o}{=}\PY{n}{X\PYZus{}test}\PY{p}{)}
\end{Verbatim}
\end{tcolorbox}

    Notice that we have trained and tested our model on two completely
separate data sets: training was performed using only the dates before
2005, and testing was performed using only the dates in 2005.

Finally, we compare the predictions for 2005 to the actual movements of
the market over that time period. We will first store the test and
training labels (recall \texttt{y\_test} is binary).

    \begin{tcolorbox}[breakable, size=fbox, boxrule=1pt, pad at break*=1mm,colback=cellbackground, colframe=cellborder]
\prompt{In}{incolor}{16}{\boxspacing}
\begin{Verbatim}[commandchars=\\\{\}]
\PY{n}{D} \PY{o}{=} \PY{n}{Smarket}\PY{o}{.}\PY{n}{Direction}
\PY{n}{L\PYZus{}train}\PY{p}{,} \PY{n}{L\PYZus{}test} \PY{o}{=} \PY{n}{D}\PY{o}{.}\PY{n}{loc}\PY{p}{[}\PY{n}{train}\PY{p}{]}\PY{p}{,} \PY{n}{D}\PY{o}{.}\PY{n}{loc}\PY{p}{[}\PY{o}{\PYZti{}}\PY{n}{train}\PY{p}{]}
\end{Verbatim}
\end{tcolorbox}

    Now we threshold the fitted probability at 50\% to form our predicted
labels.

    \begin{tcolorbox}[breakable, size=fbox, boxrule=1pt, pad at break*=1mm,colback=cellbackground, colframe=cellborder]
\prompt{In}{incolor}{17}{\boxspacing}
\begin{Verbatim}[commandchars=\\\{\}]
\PY{n}{labels} \PY{o}{=} \PY{n}{np}\PY{o}{.}\PY{n}{array}\PY{p}{(}\PY{p}{[}\PY{l+s+s1}{\PYZsq{}}\PY{l+s+s1}{Down}\PY{l+s+s1}{\PYZsq{}}\PY{p}{]}\PY{o}{*}\PY{l+m+mi}{252}\PY{p}{)}
\PY{n}{labels}\PY{p}{[}\PY{n}{probs}\PY{o}{\PYZgt{}}\PY{l+m+mf}{0.5}\PY{p}{]} \PY{o}{=} \PY{l+s+s1}{\PYZsq{}}\PY{l+s+s1}{Up}\PY{l+s+s1}{\PYZsq{}}
\PY{n}{confusion\PYZus{}table}\PY{p}{(}\PY{n}{labels}\PY{p}{,} \PY{n}{L\PYZus{}test}\PY{p}{)}
\end{Verbatim}
\end{tcolorbox}

            \begin{tcolorbox}[breakable, size=fbox, boxrule=.5pt, pad at break*=1mm, opacityfill=0]
\prompt{Out}{outcolor}{17}{\boxspacing}
\begin{Verbatim}[commandchars=\\\{\}]
Truth      Down  Up
Predicted
Down         77  97
Up           34  44
\end{Verbatim}
\end{tcolorbox}
        
    The test accuracy is about 48\% while the error rate is about 52\%

    \begin{tcolorbox}[breakable, size=fbox, boxrule=1pt, pad at break*=1mm,colback=cellbackground, colframe=cellborder]
\prompt{In}{incolor}{18}{\boxspacing}
\begin{Verbatim}[commandchars=\\\{\}]
\PY{n}{np}\PY{o}{.}\PY{n}{mean}\PY{p}{(}\PY{n}{labels} \PY{o}{==} \PY{n}{L\PYZus{}test}\PY{p}{)}\PY{p}{,} \PY{n}{np}\PY{o}{.}\PY{n}{mean}\PY{p}{(}\PY{n}{labels} \PY{o}{!=} \PY{n}{L\PYZus{}test}\PY{p}{)}
\end{Verbatim}
\end{tcolorbox}

            \begin{tcolorbox}[breakable, size=fbox, boxrule=.5pt, pad at break*=1mm, opacityfill=0]
\prompt{Out}{outcolor}{18}{\boxspacing}
\begin{Verbatim}[commandchars=\\\{\}]
(0.4801587301587302, 0.5198412698412699)
\end{Verbatim}
\end{tcolorbox}
        
    The \texttt{!=} notation means \emph{not equal to}, and so the last
command computes the test set error rate. The results are rather
disappointing: the test error rate is 52\%, which is worse than random
guessing! Of course this result is not all that surprising, given that
one would not generally expect to be able to use previous days' returns
to predict future market performance. (After all, if it were possible to
do so, then the authors of this book would be out striking it rich
rather than writing a statistics textbook.)

We recall that the logistic regression model had very underwhelming
\emph{p}-values associated with all of the predictors, and that the
smallest \emph{p}-value, though not very small, corresponded to
\texttt{Lag1}. Perhaps by removing the variables that appear not to be
helpful in predicting \texttt{Direction}, we can obtain a more effective
model. After all, using predictors that have no relationship with the
response tends to cause a deterioration in the test error rate (since
such predictors cause an increase in variance without a corresponding
decrease in bias), and so removing such predictors may in turn yield an
improvement. Below we refit the logistic regression using just
\texttt{Lag1} and \texttt{Lag2}, which seemed to have the highest
predictive power in the original logistic regression model.

    \begin{tcolorbox}[breakable, size=fbox, boxrule=1pt, pad at break*=1mm,colback=cellbackground, colframe=cellborder]
\prompt{In}{incolor}{19}{\boxspacing}
\begin{Verbatim}[commandchars=\\\{\}]
\PY{n}{model} \PY{o}{=} \PY{n}{MS}\PY{p}{(}\PY{p}{[}\PY{l+s+s1}{\PYZsq{}}\PY{l+s+s1}{Lag1}\PY{l+s+s1}{\PYZsq{}}\PY{p}{,} \PY{l+s+s1}{\PYZsq{}}\PY{l+s+s1}{Lag2}\PY{l+s+s1}{\PYZsq{}}\PY{p}{]}\PY{p}{)}\PY{o}{.}\PY{n}{fit}\PY{p}{(}\PY{n}{Smarket}\PY{p}{)}
\PY{n}{X} \PY{o}{=} \PY{n}{model}\PY{o}{.}\PY{n}{transform}\PY{p}{(}\PY{n}{Smarket}\PY{p}{)}
\PY{n}{X\PYZus{}train}\PY{p}{,} \PY{n}{X\PYZus{}test} \PY{o}{=} \PY{n}{X}\PY{o}{.}\PY{n}{loc}\PY{p}{[}\PY{n}{train}\PY{p}{]}\PY{p}{,} \PY{n}{X}\PY{o}{.}\PY{n}{loc}\PY{p}{[}\PY{o}{\PYZti{}}\PY{n}{train}\PY{p}{]}
\PY{n}{glm\PYZus{}train} \PY{o}{=} \PY{n}{sm}\PY{o}{.}\PY{n}{GLM}\PY{p}{(}\PY{n}{y\PYZus{}train}\PY{p}{,}
                   \PY{n}{X\PYZus{}train}\PY{p}{,}
                   \PY{n}{family}\PY{o}{=}\PY{n}{sm}\PY{o}{.}\PY{n}{families}\PY{o}{.}\PY{n}{Binomial}\PY{p}{(}\PY{p}{)}\PY{p}{)}
\PY{n}{results} \PY{o}{=} \PY{n}{glm\PYZus{}train}\PY{o}{.}\PY{n}{fit}\PY{p}{(}\PY{p}{)}
\PY{n}{probs} \PY{o}{=} \PY{n}{results}\PY{o}{.}\PY{n}{predict}\PY{p}{(}\PY{n}{exog}\PY{o}{=}\PY{n}{X\PYZus{}test}\PY{p}{)}
\PY{n}{labels} \PY{o}{=} \PY{n}{np}\PY{o}{.}\PY{n}{array}\PY{p}{(}\PY{p}{[}\PY{l+s+s1}{\PYZsq{}}\PY{l+s+s1}{Down}\PY{l+s+s1}{\PYZsq{}}\PY{p}{]}\PY{o}{*}\PY{l+m+mi}{252}\PY{p}{)}
\PY{n}{labels}\PY{p}{[}\PY{n}{probs}\PY{o}{\PYZgt{}}\PY{l+m+mf}{0.5}\PY{p}{]} \PY{o}{=} \PY{l+s+s1}{\PYZsq{}}\PY{l+s+s1}{Up}\PY{l+s+s1}{\PYZsq{}}
\PY{n}{confusion\PYZus{}table}\PY{p}{(}\PY{n}{labels}\PY{p}{,} \PY{n}{L\PYZus{}test}\PY{p}{)}
\end{Verbatim}
\end{tcolorbox}

            \begin{tcolorbox}[breakable, size=fbox, boxrule=.5pt, pad at break*=1mm, opacityfill=0]
\prompt{Out}{outcolor}{19}{\boxspacing}
\begin{Verbatim}[commandchars=\\\{\}]
Truth      Down   Up
Predicted
Down         35   35
Up           76  106
\end{Verbatim}
\end{tcolorbox}
        
    Let's evaluate the overall accuracy as well as the accuracy within the
days when logistic regression predicts an increase.

    \begin{tcolorbox}[breakable, size=fbox, boxrule=1pt, pad at break*=1mm,colback=cellbackground, colframe=cellborder]
\prompt{In}{incolor}{20}{\boxspacing}
\begin{Verbatim}[commandchars=\\\{\}]
\PY{p}{(}\PY{l+m+mi}{35}\PY{o}{+}\PY{l+m+mi}{106}\PY{p}{)}\PY{o}{/}\PY{l+m+mi}{252}\PY{p}{,}\PY{l+m+mi}{106}\PY{o}{/}\PY{p}{(}\PY{l+m+mi}{106}\PY{o}{+}\PY{l+m+mi}{76}\PY{p}{)}
\end{Verbatim}
\end{tcolorbox}

            \begin{tcolorbox}[breakable, size=fbox, boxrule=.5pt, pad at break*=1mm, opacityfill=0]
\prompt{Out}{outcolor}{20}{\boxspacing}
\begin{Verbatim}[commandchars=\\\{\}]
(0.5595238095238095, 0.5824175824175825)
\end{Verbatim}
\end{tcolorbox}
        
    Now the results appear to be a little better: 56\% of the daily
movements have been correctly predicted. It is worth noting that in this
case, a much simpler strategy of predicting that the market will
increase every day will also be correct 56\% of the time! Hence, in
terms of overall error rate, the logistic regression method is no better
than the naive approach. However, the confusion matrix shows that on
days when logistic regression predicts an increase in the market, it has
a 58\% accuracy rate. This suggests a possible trading strategy of
buying on days when the model predicts an increasing market, and
avoiding trades on days when a decrease is predicted. Of course one
would need to investigate more carefully whether this small improvement
was real or just due to random chance.

Suppose that we want to predict the returns associated with particular
values of \texttt{Lag1} and \texttt{Lag2}. In particular, we want to
predict \texttt{Direction} on a day when \texttt{Lag1} and \texttt{Lag2}
equal \(1.2\) and \(1.1\), respectively, and on a day when they equal
\(1.5\) and \(-0.8\). We do this using the \texttt{predict()} function.

    \begin{tcolorbox}[breakable, size=fbox, boxrule=1pt, pad at break*=1mm,colback=cellbackground, colframe=cellborder]
\prompt{In}{incolor}{21}{\boxspacing}
\begin{Verbatim}[commandchars=\\\{\}]
\PY{n}{newdata} \PY{o}{=} \PY{n}{pd}\PY{o}{.}\PY{n}{DataFrame}\PY{p}{(}\PY{p}{\PYZob{}}\PY{l+s+s1}{\PYZsq{}}\PY{l+s+s1}{Lag1}\PY{l+s+s1}{\PYZsq{}}\PY{p}{:}\PY{p}{[}\PY{l+m+mf}{1.2}\PY{p}{,} \PY{l+m+mf}{1.5}\PY{p}{]}\PY{p}{,}
                        \PY{l+s+s1}{\PYZsq{}}\PY{l+s+s1}{Lag2}\PY{l+s+s1}{\PYZsq{}}\PY{p}{:}\PY{p}{[}\PY{l+m+mf}{1.1}\PY{p}{,} \PY{o}{\PYZhy{}}\PY{l+m+mf}{0.8}\PY{p}{]}\PY{p}{\PYZcb{}}\PY{p}{)}\PY{p}{;}
\PY{n}{newX} \PY{o}{=} \PY{n}{model}\PY{o}{.}\PY{n}{transform}\PY{p}{(}\PY{n}{newdata}\PY{p}{)}
\PY{n}{results}\PY{o}{.}\PY{n}{predict}\PY{p}{(}\PY{n}{newX}\PY{p}{)}
\end{Verbatim}
\end{tcolorbox}

            \begin{tcolorbox}[breakable, size=fbox, boxrule=.5pt, pad at break*=1mm, opacityfill=0]
\prompt{Out}{outcolor}{21}{\boxspacing}
\begin{Verbatim}[commandchars=\\\{\}]
0    0.479146
1    0.496094
dtype: float64
\end{Verbatim}
\end{tcolorbox}
        
    \hypertarget{linear-discriminant-analysis}{%
\subsection{Linear Discriminant
Analysis}\label{linear-discriminant-analysis}}

    We begin by performing LDA on the \texttt{Smarket} data, using the
function \texttt{LinearDiscriminantAnalysis()}, which we have
abbreviated \texttt{LDA()}. We fit the model using only the observations
before 2005.

    \begin{tcolorbox}[breakable, size=fbox, boxrule=1pt, pad at break*=1mm,colback=cellbackground, colframe=cellborder]
\prompt{In}{incolor}{22}{\boxspacing}
\begin{Verbatim}[commandchars=\\\{\}]
\PY{n}{lda} \PY{o}{=} \PY{n}{LDA}\PY{p}{(}\PY{n}{store\PYZus{}covariance}\PY{o}{=}\PY{k+kc}{True}\PY{p}{)}
\end{Verbatim}
\end{tcolorbox}

    Since the \texttt{LDA} estimator automatically adds an intercept, we
should remove the column corresponding to the intercept in both
\texttt{X\_train} and \texttt{X\_test}. We can also directly use the
labels rather than the Boolean vectors \texttt{y\_train}.

    \begin{tcolorbox}[breakable, size=fbox, boxrule=1pt, pad at break*=1mm,colback=cellbackground, colframe=cellborder]
\prompt{In}{incolor}{23}{\boxspacing}
\begin{Verbatim}[commandchars=\\\{\}]
\PY{n}{X\PYZus{}train}\PY{p}{,} \PY{n}{X\PYZus{}test} \PY{o}{=} \PY{p}{[}\PY{n}{M}\PY{o}{.}\PY{n}{drop}\PY{p}{(}\PY{n}{columns}\PY{o}{=}\PY{p}{[}\PY{l+s+s1}{\PYZsq{}}\PY{l+s+s1}{intercept}\PY{l+s+s1}{\PYZsq{}}\PY{p}{]}\PY{p}{)}
                   \PY{k}{for} \PY{n}{M} \PY{o+ow}{in} \PY{p}{[}\PY{n}{X\PYZus{}train}\PY{p}{,} \PY{n}{X\PYZus{}test}\PY{p}{]}\PY{p}{]}
\PY{n}{lda}\PY{o}{.}\PY{n}{fit}\PY{p}{(}\PY{n}{X\PYZus{}train}\PY{p}{,} \PY{n}{L\PYZus{}train}\PY{p}{)}
\end{Verbatim}
\end{tcolorbox}

            \begin{tcolorbox}[breakable, size=fbox, boxrule=.5pt, pad at break*=1mm, opacityfill=0]
\prompt{Out}{outcolor}{23}{\boxspacing}
\begin{Verbatim}[commandchars=\\\{\}]
LinearDiscriminantAnalysis(store\_covariance=True)
\end{Verbatim}
\end{tcolorbox}
        
    Here we have used the list comprehensions introduced in Section 3.6.4.
Looking at our first line above, we see that the right-hand side is a
list of length two. This is because the code
\texttt{for\ M\ in\ {[}X\_train,\ X\_test{]}} iterates over a list of
length two. While here we loop over a list, the list comprehension
method works when looping over any iterable object. We then apply the
\texttt{drop()} method to each element in the iteration, collecting the
result in a list. The left-hand side tells \texttt{Python} to unpack
this list of length two, assigning its elements to the variables
\texttt{X\_train} and \texttt{X\_test}. Of course, this overwrites the
previous values of \texttt{X\_train} and \texttt{X\_test}.

Having fit the model, we can extract the means in the two classes with
the \texttt{means\_} attribute. These are the average of each predictor
within each class, and are used by LDA as estimates of \(\mu_k\). These
suggest that there is a tendency for the previous 2 days' returns to be
negative on days when the market increases, and a tendency for the
previous days' returns to be positive on days when the market declines.

    \begin{tcolorbox}[breakable, size=fbox, boxrule=1pt, pad at break*=1mm,colback=cellbackground, colframe=cellborder]
\prompt{In}{incolor}{24}{\boxspacing}
\begin{Verbatim}[commandchars=\\\{\}]
\PY{n}{lda}\PY{o}{.}\PY{n}{means\PYZus{}}
\end{Verbatim}
\end{tcolorbox}

            \begin{tcolorbox}[breakable, size=fbox, boxrule=.5pt, pad at break*=1mm, opacityfill=0]
\prompt{Out}{outcolor}{24}{\boxspacing}
\begin{Verbatim}[commandchars=\\\{\}]
array([[ 0.04279022,  0.03389409],
       [-0.03954635, -0.03132544]])
\end{Verbatim}
\end{tcolorbox}
        
    The estimated prior probabilities are stored in the \texttt{priors\_}
attribute. The package \texttt{sklearn} typically uses this trailing
\texttt{\_} to denote a quantity estimated when using the \texttt{fit()}
method. We can be sure of which entry corresponds to which label by
looking at the \texttt{classes\_} attribute.

    \begin{tcolorbox}[breakable, size=fbox, boxrule=1pt, pad at break*=1mm,colback=cellbackground, colframe=cellborder]
\prompt{In}{incolor}{25}{\boxspacing}
\begin{Verbatim}[commandchars=\\\{\}]
\PY{n}{lda}\PY{o}{.}\PY{n}{classes\PYZus{}}
\end{Verbatim}
\end{tcolorbox}

            \begin{tcolorbox}[breakable, size=fbox, boxrule=.5pt, pad at break*=1mm, opacityfill=0]
\prompt{Out}{outcolor}{25}{\boxspacing}
\begin{Verbatim}[commandchars=\\\{\}]
array(['Down', 'Up'], dtype='<U4')
\end{Verbatim}
\end{tcolorbox}
        
    The LDA output indicates that \(\hat\pi_{Down}=0.492\) and
\(\hat\pi_{Up}=0.508\).

    \begin{tcolorbox}[breakable, size=fbox, boxrule=1pt, pad at break*=1mm,colback=cellbackground, colframe=cellborder]
\prompt{In}{incolor}{26}{\boxspacing}
\begin{Verbatim}[commandchars=\\\{\}]
\PY{n}{lda}\PY{o}{.}\PY{n}{priors\PYZus{}}
\end{Verbatim}
\end{tcolorbox}

            \begin{tcolorbox}[breakable, size=fbox, boxrule=.5pt, pad at break*=1mm, opacityfill=0]
\prompt{Out}{outcolor}{26}{\boxspacing}
\begin{Verbatim}[commandchars=\\\{\}]
array([0.49198397, 0.50801603])
\end{Verbatim}
\end{tcolorbox}
        
    The linear discriminant vectors can be found in the \texttt{scalings\_}
attribute:

    \begin{tcolorbox}[breakable, size=fbox, boxrule=1pt, pad at break*=1mm,colback=cellbackground, colframe=cellborder]
\prompt{In}{incolor}{27}{\boxspacing}
\begin{Verbatim}[commandchars=\\\{\}]
\PY{n}{lda}\PY{o}{.}\PY{n}{scalings\PYZus{}}
\end{Verbatim}
\end{tcolorbox}

            \begin{tcolorbox}[breakable, size=fbox, boxrule=.5pt, pad at break*=1mm, opacityfill=0]
\prompt{Out}{outcolor}{27}{\boxspacing}
\begin{Verbatim}[commandchars=\\\{\}]
array([[-0.64201904],
       [-0.51352928]])
\end{Verbatim}
\end{tcolorbox}
        
    These values provide the linear combination of \texttt{Lag1} and
\texttt{Lag2} that are used to form the LDA decision rule. In other
words, these are the multipliers of the elements of \(X=x\) in (4.24).
If \$-0.64\times \texttt{Lag1} - 0.51 \times \texttt{Lag2} \$ is large,
then the LDA classifier will predict a market increase, and if it is
small, then the LDA classifier will predict a market decline.

    \begin{tcolorbox}[breakable, size=fbox, boxrule=1pt, pad at break*=1mm,colback=cellbackground, colframe=cellborder]
\prompt{In}{incolor}{28}{\boxspacing}
\begin{Verbatim}[commandchars=\\\{\}]
\PY{n}{lda\PYZus{}pred} \PY{o}{=} \PY{n}{lda}\PY{o}{.}\PY{n}{predict}\PY{p}{(}\PY{n}{X\PYZus{}test}\PY{p}{)}
\end{Verbatim}
\end{tcolorbox}

    As we observed in our comparison of classification methods (Section
4.5), the LDA and logistic regression predictions are almost identical.

    \begin{tcolorbox}[breakable, size=fbox, boxrule=1pt, pad at break*=1mm,colback=cellbackground, colframe=cellborder]
\prompt{In}{incolor}{29}{\boxspacing}
\begin{Verbatim}[commandchars=\\\{\}]
\PY{n}{confusion\PYZus{}table}\PY{p}{(}\PY{n}{lda\PYZus{}pred}\PY{p}{,} \PY{n}{L\PYZus{}test}\PY{p}{)}
\end{Verbatim}
\end{tcolorbox}

            \begin{tcolorbox}[breakable, size=fbox, boxrule=.5pt, pad at break*=1mm, opacityfill=0]
\prompt{Out}{outcolor}{29}{\boxspacing}
\begin{Verbatim}[commandchars=\\\{\}]
Truth      Down   Up
Predicted
Down         35   35
Up           76  106
\end{Verbatim}
\end{tcolorbox}
        
    We can also estimate the probability of each class for each point in a
training set. Applying a 50\% threshold to the posterior probabilities
of being in class one allows us to recreate the predictions contained in
\texttt{lda\_pred}.

    \begin{tcolorbox}[breakable, size=fbox, boxrule=1pt, pad at break*=1mm,colback=cellbackground, colframe=cellborder]
\prompt{In}{incolor}{30}{\boxspacing}
\begin{Verbatim}[commandchars=\\\{\}]
\PY{n}{lda\PYZus{}prob} \PY{o}{=} \PY{n}{lda}\PY{o}{.}\PY{n}{predict\PYZus{}proba}\PY{p}{(}\PY{n}{X\PYZus{}test}\PY{p}{)}
\PY{n}{np}\PY{o}{.}\PY{n}{all}\PY{p}{(}
       \PY{n}{np}\PY{o}{.}\PY{n}{where}\PY{p}{(}\PY{n}{lda\PYZus{}prob}\PY{p}{[}\PY{p}{:}\PY{p}{,}\PY{l+m+mi}{1}\PY{p}{]} \PY{o}{\PYZgt{}}\PY{o}{=} \PY{l+m+mf}{0.5}\PY{p}{,} \PY{l+s+s1}{\PYZsq{}}\PY{l+s+s1}{Up}\PY{l+s+s1}{\PYZsq{}}\PY{p}{,}\PY{l+s+s1}{\PYZsq{}}\PY{l+s+s1}{Down}\PY{l+s+s1}{\PYZsq{}}\PY{p}{)} \PY{o}{==} \PY{n}{lda\PYZus{}pred}
       \PY{p}{)}
\end{Verbatim}
\end{tcolorbox}

            \begin{tcolorbox}[breakable, size=fbox, boxrule=.5pt, pad at break*=1mm, opacityfill=0]
\prompt{Out}{outcolor}{30}{\boxspacing}
\begin{Verbatim}[commandchars=\\\{\}]
True
\end{Verbatim}
\end{tcolorbox}
        
    Above, we used the \texttt{np.where()} function that creates an array
with value \texttt{\textquotesingle{}Up\textquotesingle{}} for indices
where the second column of \texttt{lda\_prob} (the estimated posterior
probability of \texttt{\textquotesingle{}Up\textquotesingle{}}) is
greater than 0.5. For problems with more than two classes the labels are
chosen as the class whose posterior probability is highest:

    \begin{tcolorbox}[breakable, size=fbox, boxrule=1pt, pad at break*=1mm,colback=cellbackground, colframe=cellborder]
\prompt{In}{incolor}{31}{\boxspacing}
\begin{Verbatim}[commandchars=\\\{\}]
\PY{n}{np}\PY{o}{.}\PY{n}{all}\PY{p}{(}
       \PY{p}{[}\PY{n}{lda}\PY{o}{.}\PY{n}{classes\PYZus{}}\PY{p}{[}\PY{n}{i}\PY{p}{]} \PY{k}{for} \PY{n}{i} \PY{o+ow}{in} \PY{n}{np}\PY{o}{.}\PY{n}{argmax}\PY{p}{(}\PY{n}{lda\PYZus{}prob}\PY{p}{,} \PY{l+m+mi}{1}\PY{p}{)}\PY{p}{]} \PY{o}{==} \PY{n}{lda\PYZus{}pred}
       \PY{p}{)}
\end{Verbatim}
\end{tcolorbox}

            \begin{tcolorbox}[breakable, size=fbox, boxrule=.5pt, pad at break*=1mm, opacityfill=0]
\prompt{Out}{outcolor}{31}{\boxspacing}
\begin{Verbatim}[commandchars=\\\{\}]
True
\end{Verbatim}
\end{tcolorbox}
        
    If we wanted to use a posterior probability threshold other than 50\% in
order to make predictions, then we could easily do so. For instance,
suppose that we wish to predict a market decrease only if we are very
certain that the market will indeed decrease on that day --- say, if the
posterior probability is at least 90\%. We know that the first column of
\texttt{lda\_prob} corresponds to the label \texttt{Down} after having
checked the \texttt{classes\_} attribute, hence we use the column index
0 rather than 1 as we did above.

    \begin{tcolorbox}[breakable, size=fbox, boxrule=1pt, pad at break*=1mm,colback=cellbackground, colframe=cellborder]
\prompt{In}{incolor}{32}{\boxspacing}
\begin{Verbatim}[commandchars=\\\{\}]
\PY{n}{np}\PY{o}{.}\PY{n}{sum}\PY{p}{(}\PY{n}{lda\PYZus{}prob}\PY{p}{[}\PY{p}{:}\PY{p}{,}\PY{l+m+mi}{0}\PY{p}{]} \PY{o}{\PYZgt{}} \PY{l+m+mf}{0.9}\PY{p}{)}
\end{Verbatim}
\end{tcolorbox}

            \begin{tcolorbox}[breakable, size=fbox, boxrule=.5pt, pad at break*=1mm, opacityfill=0]
\prompt{Out}{outcolor}{32}{\boxspacing}
\begin{Verbatim}[commandchars=\\\{\}]
0
\end{Verbatim}
\end{tcolorbox}
        
    No days in 2005 meet that threshold! In fact, the greatest posterior
probability of decrease in all of 2005 was 52.02\%.

The LDA classifier above is the first classifier from the
\texttt{sklearn} library. We will use several other objects from this
library. The objects follow a common structure that simplifies tasks
such as cross-validation, which we will see in Chapter 5. Specifically,
the methods first create a generic classifier without referring to any
data. This classifier is then fit to data with the \texttt{fit()} method
and predictions are always produced with the \texttt{predict()} method.
This pattern of first instantiating the classifier, followed by fitting
it, and then producing predictions is an explicit design choice of
\texttt{sklearn}. This uniformity makes it possible to cleanly copy the
classifier so that it can be fit on different data; e.g.~different
training sets arising in cross-validation. This standard pattern also
allows for a predictable formation of workflows.

    \hypertarget{quadratic-discriminant-analysis}{%
\subsection{Quadratic Discriminant
Analysis}\label{quadratic-discriminant-analysis}}

We will now fit a QDA model to the \texttt{Smarket} data. QDA is
implemented via \texttt{QuadraticDiscriminantAnalysis()} in the
\texttt{sklearn} package, which we abbreviate to \texttt{QDA()}. The
syntax is very similar to \texttt{LDA()}.

    \begin{tcolorbox}[breakable, size=fbox, boxrule=1pt, pad at break*=1mm,colback=cellbackground, colframe=cellborder]
\prompt{In}{incolor}{33}{\boxspacing}
\begin{Verbatim}[commandchars=\\\{\}]
\PY{n}{qda} \PY{o}{=} \PY{n}{QDA}\PY{p}{(}\PY{n}{store\PYZus{}covariance}\PY{o}{=}\PY{k+kc}{True}\PY{p}{)}
\PY{n}{qda}\PY{o}{.}\PY{n}{fit}\PY{p}{(}\PY{n}{X\PYZus{}train}\PY{p}{,} \PY{n}{L\PYZus{}train}\PY{p}{)}
\end{Verbatim}
\end{tcolorbox}

            \begin{tcolorbox}[breakable, size=fbox, boxrule=.5pt, pad at break*=1mm, opacityfill=0]
\prompt{Out}{outcolor}{33}{\boxspacing}
\begin{Verbatim}[commandchars=\\\{\}]
QuadraticDiscriminantAnalysis(store\_covariance=True)
\end{Verbatim}
\end{tcolorbox}
        
    The \texttt{QDA()} function will again compute \texttt{means\_} and
\texttt{priors\_}.

    \begin{tcolorbox}[breakable, size=fbox, boxrule=1pt, pad at break*=1mm,colback=cellbackground, colframe=cellborder]
\prompt{In}{incolor}{34}{\boxspacing}
\begin{Verbatim}[commandchars=\\\{\}]
\PY{n}{qda}\PY{o}{.}\PY{n}{means\PYZus{}}\PY{p}{,} \PY{n}{qda}\PY{o}{.}\PY{n}{priors\PYZus{}}
\end{Verbatim}
\end{tcolorbox}

            \begin{tcolorbox}[breakable, size=fbox, boxrule=.5pt, pad at break*=1mm, opacityfill=0]
\prompt{Out}{outcolor}{34}{\boxspacing}
\begin{Verbatim}[commandchars=\\\{\}]
(array([[ 0.04279022,  0.03389409],
        [-0.03954635, -0.03132544]]),
 array([0.49198397, 0.50801603]))
\end{Verbatim}
\end{tcolorbox}
        
    The \texttt{QDA()} classifier will estimate one covariance per class.
Here is the estimated covariance in the first class:

    \begin{tcolorbox}[breakable, size=fbox, boxrule=1pt, pad at break*=1mm,colback=cellbackground, colframe=cellborder]
\prompt{In}{incolor}{35}{\boxspacing}
\begin{Verbatim}[commandchars=\\\{\}]
\PY{n}{qda}\PY{o}{.}\PY{n}{covariance\PYZus{}}\PY{p}{[}\PY{l+m+mi}{0}\PY{p}{]}
\end{Verbatim}
\end{tcolorbox}

            \begin{tcolorbox}[breakable, size=fbox, boxrule=.5pt, pad at break*=1mm, opacityfill=0]
\prompt{Out}{outcolor}{35}{\boxspacing}
\begin{Verbatim}[commandchars=\\\{\}]
array([[ 1.50662277, -0.03924806],
       [-0.03924806,  1.53559498]])
\end{Verbatim}
\end{tcolorbox}
        
    The output contains the group means. But it does not contain the
coefficients of the linear discriminants, because the QDA classifier
involves a quadratic, rather than a linear, function of the predictors.
The \texttt{predict()} function works in exactly the same fashion as for
LDA.

    \begin{tcolorbox}[breakable, size=fbox, boxrule=1pt, pad at break*=1mm,colback=cellbackground, colframe=cellborder]
\prompt{In}{incolor}{36}{\boxspacing}
\begin{Verbatim}[commandchars=\\\{\}]
\PY{n}{qda\PYZus{}pred} \PY{o}{=} \PY{n}{qda}\PY{o}{.}\PY{n}{predict}\PY{p}{(}\PY{n}{X\PYZus{}test}\PY{p}{)}
\PY{n}{confusion\PYZus{}table}\PY{p}{(}\PY{n}{qda\PYZus{}pred}\PY{p}{,} \PY{n}{L\PYZus{}test}\PY{p}{)}
\end{Verbatim}
\end{tcolorbox}

            \begin{tcolorbox}[breakable, size=fbox, boxrule=.5pt, pad at break*=1mm, opacityfill=0]
\prompt{Out}{outcolor}{36}{\boxspacing}
\begin{Verbatim}[commandchars=\\\{\}]
Truth      Down   Up
Predicted
Down         30   20
Up           81  121
\end{Verbatim}
\end{tcolorbox}
        
    Interestingly, the QDA predictions are accurate almost 60\% of the time,
even though the 2005 data was not used to fit the model.

    \begin{tcolorbox}[breakable, size=fbox, boxrule=1pt, pad at break*=1mm,colback=cellbackground, colframe=cellborder]
\prompt{In}{incolor}{37}{\boxspacing}
\begin{Verbatim}[commandchars=\\\{\}]
\PY{n}{np}\PY{o}{.}\PY{n}{mean}\PY{p}{(}\PY{n}{qda\PYZus{}pred} \PY{o}{==} \PY{n}{L\PYZus{}test}\PY{p}{)}
\end{Verbatim}
\end{tcolorbox}

            \begin{tcolorbox}[breakable, size=fbox, boxrule=.5pt, pad at break*=1mm, opacityfill=0]
\prompt{Out}{outcolor}{37}{\boxspacing}
\begin{Verbatim}[commandchars=\\\{\}]
0.5992063492063492
\end{Verbatim}
\end{tcolorbox}
        
    This level of accuracy is quite impressive for stock market data, which
is known to be quite hard to model accurately. This suggests that the
quadratic form assumed by QDA may capture the true relationship more
accurately than the linear forms assumed by LDA and logistic regression.
However, we recommend evaluating this method's performance on a larger
test set before betting that this approach will consistently beat the
market!

    \hypertarget{naive-bayes}{%
\subsection{Naive Bayes}\label{naive-bayes}}

Next, we fit a naive Bayes model to the \texttt{Smarket} data. The
syntax is similar to that of \texttt{LDA()} and \texttt{QDA()}. By
default, this implementation \texttt{GaussianNB()} of the naive Bayes
classifier models each quantitative feature using a Gaussian
distribution. However, a kernel density method can also be used to
estimate the distributions.

    \begin{tcolorbox}[breakable, size=fbox, boxrule=1pt, pad at break*=1mm,colback=cellbackground, colframe=cellborder]
\prompt{In}{incolor}{38}{\boxspacing}
\begin{Verbatim}[commandchars=\\\{\}]
\PY{n}{NB} \PY{o}{=} \PY{n}{GaussianNB}\PY{p}{(}\PY{p}{)}
\PY{n}{NB}\PY{o}{.}\PY{n}{fit}\PY{p}{(}\PY{n}{X\PYZus{}train}\PY{p}{,} \PY{n}{L\PYZus{}train}\PY{p}{)}
\end{Verbatim}
\end{tcolorbox}

            \begin{tcolorbox}[breakable, size=fbox, boxrule=.5pt, pad at break*=1mm, opacityfill=0]
\prompt{Out}{outcolor}{38}{\boxspacing}
\begin{Verbatim}[commandchars=\\\{\}]
GaussianNB()
\end{Verbatim}
\end{tcolorbox}
        
    The classes are stored as \texttt{classes\_}.

    \begin{tcolorbox}[breakable, size=fbox, boxrule=1pt, pad at break*=1mm,colback=cellbackground, colframe=cellborder]
\prompt{In}{incolor}{39}{\boxspacing}
\begin{Verbatim}[commandchars=\\\{\}]
\PY{n}{NB}\PY{o}{.}\PY{n}{classes\PYZus{}}
\end{Verbatim}
\end{tcolorbox}

            \begin{tcolorbox}[breakable, size=fbox, boxrule=.5pt, pad at break*=1mm, opacityfill=0]
\prompt{Out}{outcolor}{39}{\boxspacing}
\begin{Verbatim}[commandchars=\\\{\}]
array(['Down', 'Up'], dtype='<U4')
\end{Verbatim}
\end{tcolorbox}
        
    The class prior probabilities are stored in the \texttt{class\_prior\_}
attribute.

    \begin{tcolorbox}[breakable, size=fbox, boxrule=1pt, pad at break*=1mm,colback=cellbackground, colframe=cellborder]
\prompt{In}{incolor}{40}{\boxspacing}
\begin{Verbatim}[commandchars=\\\{\}]
\PY{n}{NB}\PY{o}{.}\PY{n}{class\PYZus{}prior\PYZus{}}
\end{Verbatim}
\end{tcolorbox}

            \begin{tcolorbox}[breakable, size=fbox, boxrule=.5pt, pad at break*=1mm, opacityfill=0]
\prompt{Out}{outcolor}{40}{\boxspacing}
\begin{Verbatim}[commandchars=\\\{\}]
array([0.49198397, 0.50801603])
\end{Verbatim}
\end{tcolorbox}
        
    The parameters of the features can be found in the \texttt{theta\_} and
\texttt{var\_} attributes. The number of rows is equal to the number of
classes, while the number of columns is equal to the number of features.
We see below that the mean for feature \texttt{Lag1} in the
\texttt{Down} class is 0.043.

    \begin{tcolorbox}[breakable, size=fbox, boxrule=1pt, pad at break*=1mm,colback=cellbackground, colframe=cellborder]
\prompt{In}{incolor}{41}{\boxspacing}
\begin{Verbatim}[commandchars=\\\{\}]
\PY{n}{NB}\PY{o}{.}\PY{n}{theta\PYZus{}}
\end{Verbatim}
\end{tcolorbox}

            \begin{tcolorbox}[breakable, size=fbox, boxrule=.5pt, pad at break*=1mm, opacityfill=0]
\prompt{Out}{outcolor}{41}{\boxspacing}
\begin{Verbatim}[commandchars=\\\{\}]
array([[ 0.04279022,  0.03389409],
       [-0.03954635, -0.03132544]])
\end{Verbatim}
\end{tcolorbox}
        
    Its variance is 1.503.

    \begin{tcolorbox}[breakable, size=fbox, boxrule=1pt, pad at break*=1mm,colback=cellbackground, colframe=cellborder]
\prompt{In}{incolor}{42}{\boxspacing}
\begin{Verbatim}[commandchars=\\\{\}]
\PY{n}{NB}\PY{o}{.}\PY{n}{var\PYZus{}}
\end{Verbatim}
\end{tcolorbox}

            \begin{tcolorbox}[breakable, size=fbox, boxrule=.5pt, pad at break*=1mm, opacityfill=0]
\prompt{Out}{outcolor}{42}{\boxspacing}
\begin{Verbatim}[commandchars=\\\{\}]
array([[1.50355429, 1.53246749],
       [1.51401364, 1.48732877]])
\end{Verbatim}
\end{tcolorbox}
        
    How do we know the names of these attributes? We use \texttt{NB?} (or
\texttt{?NB}).

We can easily verify the mean computation:

    \begin{tcolorbox}[breakable, size=fbox, boxrule=1pt, pad at break*=1mm,colback=cellbackground, colframe=cellborder]
\prompt{In}{incolor}{43}{\boxspacing}
\begin{Verbatim}[commandchars=\\\{\}]
\PY{n}{X\PYZus{}train}\PY{p}{[}\PY{n}{L\PYZus{}train} \PY{o}{==} \PY{l+s+s1}{\PYZsq{}}\PY{l+s+s1}{Down}\PY{l+s+s1}{\PYZsq{}}\PY{p}{]}\PY{o}{.}\PY{n}{mean}\PY{p}{(}\PY{p}{)}
\end{Verbatim}
\end{tcolorbox}

            \begin{tcolorbox}[breakable, size=fbox, boxrule=.5pt, pad at break*=1mm, opacityfill=0]
\prompt{Out}{outcolor}{43}{\boxspacing}
\begin{Verbatim}[commandchars=\\\{\}]
Lag1    0.042790
Lag2    0.033894
dtype: float64
\end{Verbatim}
\end{tcolorbox}
        
    Similarly for the variance:

    \begin{tcolorbox}[breakable, size=fbox, boxrule=1pt, pad at break*=1mm,colback=cellbackground, colframe=cellborder]
\prompt{In}{incolor}{44}{\boxspacing}
\begin{Verbatim}[commandchars=\\\{\}]
\PY{n}{X\PYZus{}train}\PY{p}{[}\PY{n}{L\PYZus{}train} \PY{o}{==} \PY{l+s+s1}{\PYZsq{}}\PY{l+s+s1}{Down}\PY{l+s+s1}{\PYZsq{}}\PY{p}{]}\PY{o}{.}\PY{n}{var}\PY{p}{(}\PY{n}{ddof}\PY{o}{=}\PY{l+m+mi}{0}\PY{p}{)}
\end{Verbatim}
\end{tcolorbox}

            \begin{tcolorbox}[breakable, size=fbox, boxrule=.5pt, pad at break*=1mm, opacityfill=0]
\prompt{Out}{outcolor}{44}{\boxspacing}
\begin{Verbatim}[commandchars=\\\{\}]
Lag1    1.503554
Lag2    1.532467
dtype: float64
\end{Verbatim}
\end{tcolorbox}
        
    Since \texttt{NB()} is a classifier in the \texttt{sklearn} library,
making predictions uses the same syntax as for \texttt{LDA()} and
\texttt{QDA()} above.

    \begin{tcolorbox}[breakable, size=fbox, boxrule=1pt, pad at break*=1mm,colback=cellbackground, colframe=cellborder]
\prompt{In}{incolor}{45}{\boxspacing}
\begin{Verbatim}[commandchars=\\\{\}]
\PY{n}{nb\PYZus{}labels} \PY{o}{=} \PY{n}{NB}\PY{o}{.}\PY{n}{predict}\PY{p}{(}\PY{n}{X\PYZus{}test}\PY{p}{)}
\PY{n}{confusion\PYZus{}table}\PY{p}{(}\PY{n}{nb\PYZus{}labels}\PY{p}{,} \PY{n}{L\PYZus{}test}\PY{p}{)}
\end{Verbatim}
\end{tcolorbox}

            \begin{tcolorbox}[breakable, size=fbox, boxrule=.5pt, pad at break*=1mm, opacityfill=0]
\prompt{Out}{outcolor}{45}{\boxspacing}
\begin{Verbatim}[commandchars=\\\{\}]
Truth      Down   Up
Predicted
Down         29   20
Up           82  121
\end{Verbatim}
\end{tcolorbox}
        
    Naive Bayes performs well on these data, with accurate predictions over
59\% of the time. This is slightly worse than QDA, but much better than
LDA.

As for \texttt{LDA}, the \texttt{predict\_proba()} method estimates the
probability that each observation belongs to a particular class.

    \begin{tcolorbox}[breakable, size=fbox, boxrule=1pt, pad at break*=1mm,colback=cellbackground, colframe=cellborder]
\prompt{In}{incolor}{46}{\boxspacing}
\begin{Verbatim}[commandchars=\\\{\}]
\PY{n}{NB}\PY{o}{.}\PY{n}{predict\PYZus{}proba}\PY{p}{(}\PY{n}{X\PYZus{}test}\PY{p}{)}\PY{p}{[}\PY{p}{:}\PY{l+m+mi}{5}\PY{p}{]}
\end{Verbatim}
\end{tcolorbox}

            \begin{tcolorbox}[breakable, size=fbox, boxrule=.5pt, pad at break*=1mm, opacityfill=0]
\prompt{Out}{outcolor}{46}{\boxspacing}
\begin{Verbatim}[commandchars=\\\{\}]
array([[0.4873288 , 0.5126712 ],
       [0.47623584, 0.52376416],
       [0.46529531, 0.53470469],
       [0.47484469, 0.52515531],
       [0.49020587, 0.50979413]])
\end{Verbatim}
\end{tcolorbox}
        
    \hypertarget{k-nearest-neighbors}{%
\subsection{K-Nearest Neighbors}\label{k-nearest-neighbors}}

We will now perform KNN using the \texttt{KNeighborsClassifier()}
function. This function works similarly to the other model-fitting
functions that we have encountered thus far.

As is the case for LDA and QDA, we fit the classifier using the
\texttt{fit} method. New predictions are formed using the
\texttt{predict} method of the object returned by \texttt{fit()}.

    \begin{tcolorbox}[breakable, size=fbox, boxrule=1pt, pad at break*=1mm,colback=cellbackground, colframe=cellborder]
\prompt{In}{incolor}{47}{\boxspacing}
\begin{Verbatim}[commandchars=\\\{\}]
\PY{n}{knn1} \PY{o}{=} \PY{n}{KNeighborsClassifier}\PY{p}{(}\PY{n}{n\PYZus{}neighbors}\PY{o}{=}\PY{l+m+mi}{1}\PY{p}{)}
\PY{n}{X\PYZus{}train}\PY{p}{,} \PY{n}{X\PYZus{}test} \PY{o}{=} \PY{p}{[}\PY{n}{np}\PY{o}{.}\PY{n}{asarray}\PY{p}{(}\PY{n}{X}\PY{p}{)} \PY{k}{for} \PY{n}{X} \PY{o+ow}{in} \PY{p}{[}\PY{n}{X\PYZus{}train}\PY{p}{,} \PY{n}{X\PYZus{}test}\PY{p}{]}\PY{p}{]}
\PY{n}{knn1}\PY{o}{.}\PY{n}{fit}\PY{p}{(}\PY{n}{X\PYZus{}train}\PY{p}{,} \PY{n}{L\PYZus{}train}\PY{p}{)}
\PY{n}{knn1\PYZus{}pred} \PY{o}{=} \PY{n}{knn1}\PY{o}{.}\PY{n}{predict}\PY{p}{(}\PY{n}{X\PYZus{}test}\PY{p}{)}
\PY{n}{confusion\PYZus{}table}\PY{p}{(}\PY{n}{knn1\PYZus{}pred}\PY{p}{,} \PY{n}{L\PYZus{}test}\PY{p}{)}
\end{Verbatim}
\end{tcolorbox}

            \begin{tcolorbox}[breakable, size=fbox, boxrule=.5pt, pad at break*=1mm, opacityfill=0]
\prompt{Out}{outcolor}{47}{\boxspacing}
\begin{Verbatim}[commandchars=\\\{\}]
Truth      Down  Up
Predicted
Down         43  58
Up           68  83
\end{Verbatim}
\end{tcolorbox}
        
    The results using \(K=1\) are not very good, since only \(50%
\) of the observations are correctly predicted. Of course, it may be
that \(K=1\) results in an overly-flexible fit to the data.

    \begin{tcolorbox}[breakable, size=fbox, boxrule=1pt, pad at break*=1mm,colback=cellbackground, colframe=cellborder]
\prompt{In}{incolor}{48}{\boxspacing}
\begin{Verbatim}[commandchars=\\\{\}]
\PY{p}{(}\PY{l+m+mi}{83}\PY{o}{+}\PY{l+m+mi}{43}\PY{p}{)}\PY{o}{/}\PY{l+m+mi}{252}\PY{p}{,} \PY{n}{np}\PY{o}{.}\PY{n}{mean}\PY{p}{(}\PY{n}{knn1\PYZus{}pred} \PY{o}{==} \PY{n}{L\PYZus{}test}\PY{p}{)}
\end{Verbatim}
\end{tcolorbox}

            \begin{tcolorbox}[breakable, size=fbox, boxrule=.5pt, pad at break*=1mm, opacityfill=0]
\prompt{Out}{outcolor}{48}{\boxspacing}
\begin{Verbatim}[commandchars=\\\{\}]
(0.5, 0.5)
\end{Verbatim}
\end{tcolorbox}
        
    We repeat the analysis below using \(K=3\).

    \begin{tcolorbox}[breakable, size=fbox, boxrule=1pt, pad at break*=1mm,colback=cellbackground, colframe=cellborder]
\prompt{In}{incolor}{49}{\boxspacing}
\begin{Verbatim}[commandchars=\\\{\}]
\PY{n}{knn3} \PY{o}{=} \PY{n}{KNeighborsClassifier}\PY{p}{(}\PY{n}{n\PYZus{}neighbors}\PY{o}{=}\PY{l+m+mi}{3}\PY{p}{)}
\PY{n}{knn3\PYZus{}pred} \PY{o}{=} \PY{n}{knn3}\PY{o}{.}\PY{n}{fit}\PY{p}{(}\PY{n}{X\PYZus{}train}\PY{p}{,} \PY{n}{L\PYZus{}train}\PY{p}{)}\PY{o}{.}\PY{n}{predict}\PY{p}{(}\PY{n}{X\PYZus{}test}\PY{p}{)}
\PY{n}{np}\PY{o}{.}\PY{n}{mean}\PY{p}{(}\PY{n}{knn3\PYZus{}pred} \PY{o}{==} \PY{n}{L\PYZus{}test}\PY{p}{)}
\end{Verbatim}
\end{tcolorbox}

            \begin{tcolorbox}[breakable, size=fbox, boxrule=.5pt, pad at break*=1mm, opacityfill=0]
\prompt{Out}{outcolor}{49}{\boxspacing}
\begin{Verbatim}[commandchars=\\\{\}]
0.5317460317460317
\end{Verbatim}
\end{tcolorbox}
        
    The results have improved slightly. But increasing \emph{K} further
provides no further improvements. It appears that for these data, and
this train/test split, QDA gives the best results of the methods that we
have examined so far.

    KNN does not perform well on the \texttt{Smarket} data, but it often
does provide impressive results. As an example we will apply the KNN
approach to the \texttt{Caravan} data set, which is part of the
\texttt{ISLP} library. This data set includes 85 predictors that measure
demographic characteristics for 5,822 individuals. The response variable
is \texttt{Purchase}, which indicates whether or not a given individual
purchases a caravan insurance policy. In this data set, only 6\% of
people purchased caravan insurance.

    \begin{tcolorbox}[breakable, size=fbox, boxrule=1pt, pad at break*=1mm,colback=cellbackground, colframe=cellborder]
\prompt{In}{incolor}{50}{\boxspacing}
\begin{Verbatim}[commandchars=\\\{\}]
\PY{n}{Caravan} \PY{o}{=} \PY{n}{load\PYZus{}data}\PY{p}{(}\PY{l+s+s1}{\PYZsq{}}\PY{l+s+s1}{Caravan}\PY{l+s+s1}{\PYZsq{}}\PY{p}{)}
\PY{n}{Purchase} \PY{o}{=} \PY{n}{Caravan}\PY{o}{.}\PY{n}{Purchase}
\PY{n}{Purchase}\PY{o}{.}\PY{n}{value\PYZus{}counts}\PY{p}{(}\PY{p}{)}
\end{Verbatim}
\end{tcolorbox}

            \begin{tcolorbox}[breakable, size=fbox, boxrule=.5pt, pad at break*=1mm, opacityfill=0]
\prompt{Out}{outcolor}{50}{\boxspacing}
\begin{Verbatim}[commandchars=\\\{\}]
Purchase
No     5474
Yes     348
Name: count, dtype: int64
\end{Verbatim}
\end{tcolorbox}
        
    The method \texttt{value\_counts()} takes a \texttt{pd.Series} or
\texttt{pd.DataFrame} and returns a \texttt{pd.Series} with the
corresponding counts for each unique element. In this case
\texttt{Purchase} has only \texttt{Yes} and \texttt{No} values and the
method returns how many values of each there are.

    \begin{tcolorbox}[breakable, size=fbox, boxrule=1pt, pad at break*=1mm,colback=cellbackground, colframe=cellborder]
\prompt{In}{incolor}{51}{\boxspacing}
\begin{Verbatim}[commandchars=\\\{\}]
\PY{l+m+mi}{348} \PY{o}{/} \PY{l+m+mi}{5822}
\end{Verbatim}
\end{tcolorbox}

            \begin{tcolorbox}[breakable, size=fbox, boxrule=.5pt, pad at break*=1mm, opacityfill=0]
\prompt{Out}{outcolor}{51}{\boxspacing}
\begin{Verbatim}[commandchars=\\\{\}]
0.05977327378907592
\end{Verbatim}
\end{tcolorbox}
        
    Our features will include all columns except \texttt{Purchase}.

    \begin{tcolorbox}[breakable, size=fbox, boxrule=1pt, pad at break*=1mm,colback=cellbackground, colframe=cellborder]
\prompt{In}{incolor}{52}{\boxspacing}
\begin{Verbatim}[commandchars=\\\{\}]
\PY{n}{feature\PYZus{}df} \PY{o}{=} \PY{n}{Caravan}\PY{o}{.}\PY{n}{drop}\PY{p}{(}\PY{n}{columns}\PY{o}{=}\PY{p}{[}\PY{l+s+s1}{\PYZsq{}}\PY{l+s+s1}{Purchase}\PY{l+s+s1}{\PYZsq{}}\PY{p}{]}\PY{p}{)}
\end{Verbatim}
\end{tcolorbox}

    Because the KNN classifier predicts the class of a given test
observation by identifying the observations that are nearest to it, the
scale of the variables matters. Any variables that are on a large scale
will have a much larger effect on the \emph{distance} between the
observations, and hence on the KNN classifier, than variables that are
on a small scale. For instance, imagine a data set that contains two
variables, \texttt{salary} and \texttt{age} (measured in dollars and
years, respectively). As far as KNN is concerned, a difference of 1,000
USD in salary is enormous compared to a difference of 50 years in age.
Consequently, \texttt{salary} will drive the KNN classification results,
and \texttt{age} will have almost no effect. This is contrary to our
intuition that a salary difference of 1,000 USD is quite small compared
to an age difference of 50 years. Furthermore, the importance of scale
to the KNN classifier leads to another issue: if we measured
\texttt{salary} in Japanese yen, or if we measured \texttt{age} in
minutes, then we'd get quite different classification results from what
we get if these two variables are measured in dollars and years.

A good way to handle this problem is to \emph{standardize} the data so
that all variables are given a mean of zero and a standard deviation of
one. Then all variables will be on a comparable scale. This is
accomplished using the \texttt{StandardScaler()} transformation.

    \begin{tcolorbox}[breakable, size=fbox, boxrule=1pt, pad at break*=1mm,colback=cellbackground, colframe=cellborder]
\prompt{In}{incolor}{53}{\boxspacing}
\begin{Verbatim}[commandchars=\\\{\}]
\PY{n}{scaler} \PY{o}{=} \PY{n}{StandardScaler}\PY{p}{(}\PY{n}{with\PYZus{}mean}\PY{o}{=}\PY{k+kc}{True}\PY{p}{,}
                        \PY{n}{with\PYZus{}std}\PY{o}{=}\PY{k+kc}{True}\PY{p}{,}
                        \PY{n}{copy}\PY{o}{=}\PY{k+kc}{True}\PY{p}{)}
\end{Verbatim}
\end{tcolorbox}

    The argument \texttt{with\_mean} indicates whether or not we should
subtract the mean, while \texttt{with\_std} indicates whether or not we
should scale the columns to have standard deviation of 1 or not.
Finally, the argument \texttt{copy=True} indicates that we will always
copy data, rather than trying to do calculations in place where
possible.

This transformation can be fit and then applied to arbitrary data. In
the first line below, the parameters for the scaling are computed and
stored in \texttt{scaler}, while the second line actually constructs the
standardized set of features.

    \begin{tcolorbox}[breakable, size=fbox, boxrule=1pt, pad at break*=1mm,colback=cellbackground, colframe=cellborder]
\prompt{In}{incolor}{54}{\boxspacing}
\begin{Verbatim}[commandchars=\\\{\}]
\PY{n}{scaler}\PY{o}{.}\PY{n}{fit}\PY{p}{(}\PY{n}{feature\PYZus{}df}\PY{p}{)}
\PY{n}{X\PYZus{}std} \PY{o}{=} \PY{n}{scaler}\PY{o}{.}\PY{n}{transform}\PY{p}{(}\PY{n}{feature\PYZus{}df}\PY{p}{)}
\end{Verbatim}
\end{tcolorbox}

    Now every column of \texttt{feature\_std} below has a standard deviation
of one and a mean of zero.

    \begin{tcolorbox}[breakable, size=fbox, boxrule=1pt, pad at break*=1mm,colback=cellbackground, colframe=cellborder]
\prompt{In}{incolor}{55}{\boxspacing}
\begin{Verbatim}[commandchars=\\\{\}]
\PY{n}{feature\PYZus{}std} \PY{o}{=} \PY{n}{pd}\PY{o}{.}\PY{n}{DataFrame}\PY{p}{(}
                 \PY{n}{X\PYZus{}std}\PY{p}{,}
                 \PY{n}{columns}\PY{o}{=}\PY{n}{feature\PYZus{}df}\PY{o}{.}\PY{n}{columns}\PY{p}{)}\PY{p}{;}
\PY{n}{feature\PYZus{}std}\PY{o}{.}\PY{n}{std}\PY{p}{(}\PY{p}{)}
\end{Verbatim}
\end{tcolorbox}

            \begin{tcolorbox}[breakable, size=fbox, boxrule=.5pt, pad at break*=1mm, opacityfill=0]
\prompt{Out}{outcolor}{55}{\boxspacing}
\begin{Verbatim}[commandchars=\\\{\}]
MOSTYPE     1.000086
MAANTHUI    1.000086
MGEMOMV     1.000086
MGEMLEEF    1.000086
MOSHOOFD    1.000086
              {\ldots}
AZEILPL     1.000086
APLEZIER    1.000086
AFIETS      1.000086
AINBOED     1.000086
ABYSTAND    1.000086
Length: 85, dtype: float64
\end{Verbatim}
\end{tcolorbox}
        
    Notice that the standard deviations are not quite \(1\) here; this is
again due to some procedures using the \(1/n\) convention for variances
(in this case \texttt{scaler()}), while others use \(1/(n-1)\) (the
\texttt{std()} method). See the footnote on page 183. In this case it
does not matter, as long as the variables are all on the same scale.

Using the function \texttt{train\_test\_split()} we now split the
observations into a test set, containing 1000 observations, and a
training set containing the remaining observations. The argument
\texttt{random\_state=0} ensures that we get the same split each time we
rerun the code.

    \begin{tcolorbox}[breakable, size=fbox, boxrule=1pt, pad at break*=1mm,colback=cellbackground, colframe=cellborder]
\prompt{In}{incolor}{56}{\boxspacing}
\begin{Verbatim}[commandchars=\\\{\}]
\PY{p}{(}\PY{n}{X\PYZus{}train}\PY{p}{,}
 \PY{n}{X\PYZus{}test}\PY{p}{,}
 \PY{n}{y\PYZus{}train}\PY{p}{,}
 \PY{n}{y\PYZus{}test}\PY{p}{)} \PY{o}{=} \PY{n}{train\PYZus{}test\PYZus{}split}\PY{p}{(}\PY{n}{np}\PY{o}{.}\PY{n}{asarray}\PY{p}{(}\PY{n}{feature\PYZus{}std}\PY{p}{)}\PY{p}{,}
                            \PY{n}{Purchase}\PY{p}{,}
                            \PY{n}{test\PYZus{}size}\PY{o}{=}\PY{l+m+mi}{1000}\PY{p}{,}
                            \PY{n}{random\PYZus{}state}\PY{o}{=}\PY{l+m+mi}{0}\PY{p}{)}
\end{Verbatim}
\end{tcolorbox}

    \texttt{?train\_test\_split} reveals that the non-keyword arguments can
be \texttt{lists}, \texttt{arrays}, \texttt{pandas\ dataframes} etc that
all have the same length (\texttt{shape{[}0{]}}) and hence are
\emph{indexable}. In this case they are the dataframe
\texttt{feature\_std} and the response variable \texttt{Purchase}.
\{Note that we have converted \texttt{feature\_std} to an
\texttt{ndarray} to address a bug in \texttt{sklearn}.\} We fit a KNN
model on the training data using \(K=1\), and evaluate its performance
on the test data.

    \begin{tcolorbox}[breakable, size=fbox, boxrule=1pt, pad at break*=1mm,colback=cellbackground, colframe=cellborder]
\prompt{In}{incolor}{57}{\boxspacing}
\begin{Verbatim}[commandchars=\\\{\}]
\PY{n}{knn1} \PY{o}{=} \PY{n}{KNeighborsClassifier}\PY{p}{(}\PY{n}{n\PYZus{}neighbors}\PY{o}{=}\PY{l+m+mi}{1}\PY{p}{)}
\PY{n}{knn1\PYZus{}pred} \PY{o}{=} \PY{n}{knn1}\PY{o}{.}\PY{n}{fit}\PY{p}{(}\PY{n}{X\PYZus{}train}\PY{p}{,} \PY{n}{y\PYZus{}train}\PY{p}{)}\PY{o}{.}\PY{n}{predict}\PY{p}{(}\PY{n}{X\PYZus{}test}\PY{p}{)}
\PY{n}{np}\PY{o}{.}\PY{n}{mean}\PY{p}{(}\PY{n}{y\PYZus{}test} \PY{o}{!=} \PY{n}{knn1\PYZus{}pred}\PY{p}{)}\PY{p}{,} \PY{n}{np}\PY{o}{.}\PY{n}{mean}\PY{p}{(}\PY{n}{y\PYZus{}test} \PY{o}{!=} \PY{l+s+s2}{\PYZdq{}}\PY{l+s+s2}{No}\PY{l+s+s2}{\PYZdq{}}\PY{p}{)}
\end{Verbatim}
\end{tcolorbox}

            \begin{tcolorbox}[breakable, size=fbox, boxrule=.5pt, pad at break*=1mm, opacityfill=0]
\prompt{Out}{outcolor}{57}{\boxspacing}
\begin{Verbatim}[commandchars=\\\{\}]
(0.111, 0.067)
\end{Verbatim}
\end{tcolorbox}
        
    The KNN error rate on the 1,000 test observations is about \(11%
\). At first glance, this may appear to be fairly good. However, since
just over 6\% of customers purchased insurance, we could get the error
rate down to almost 6\% by always predicting \texttt{No} regardless of
the values of the predictors! This is known as the \emph{null rate}.\}

Suppose that there is some non-trivial cost to trying to sell insurance
to a given individual. For instance, perhaps a salesperson must visit
each potential customer. If the company tries to sell insurance to a
random selection of customers, then the success rate will be only 6\%,
which may be far too low given the costs involved. Instead, the company
would like to try to sell insurance only to customers who are likely to
buy it. So the overall error rate is not of interest. Instead, the
fraction of individuals that are correctly predicted to buy insurance is
of interest.

    \begin{tcolorbox}[breakable, size=fbox, boxrule=1pt, pad at break*=1mm,colback=cellbackground, colframe=cellborder]
\prompt{In}{incolor}{58}{\boxspacing}
\begin{Verbatim}[commandchars=\\\{\}]
\PY{n}{confusion\PYZus{}table}\PY{p}{(}\PY{n}{knn1\PYZus{}pred}\PY{p}{,} \PY{n}{y\PYZus{}test}\PY{p}{)}
\end{Verbatim}
\end{tcolorbox}

            \begin{tcolorbox}[breakable, size=fbox, boxrule=.5pt, pad at break*=1mm, opacityfill=0]
\prompt{Out}{outcolor}{58}{\boxspacing}
\begin{Verbatim}[commandchars=\\\{\}]
Truth       No  Yes
Predicted
No         880   58
Yes         53    9
\end{Verbatim}
\end{tcolorbox}
        
    It turns out that KNN with \(K=1\) does far better than random guessing
among the customers that are predicted to buy insurance. Among 62 such
customers, 9, or 14.5\%, actually do purchase insurance. This is double
the rate that one would obtain from random guessing.

    \begin{tcolorbox}[breakable, size=fbox, boxrule=1pt, pad at break*=1mm,colback=cellbackground, colframe=cellborder]
\prompt{In}{incolor}{59}{\boxspacing}
\begin{Verbatim}[commandchars=\\\{\}]
\PY{l+m+mi}{9}\PY{o}{/}\PY{p}{(}\PY{l+m+mi}{53}\PY{o}{+}\PY{l+m+mi}{9}\PY{p}{)}
\end{Verbatim}
\end{tcolorbox}

            \begin{tcolorbox}[breakable, size=fbox, boxrule=.5pt, pad at break*=1mm, opacityfill=0]
\prompt{Out}{outcolor}{59}{\boxspacing}
\begin{Verbatim}[commandchars=\\\{\}]
0.14516129032258066
\end{Verbatim}
\end{tcolorbox}
        
    \hypertarget{tuning-parameters}{%
\subsubsection{Tuning Parameters}\label{tuning-parameters}}

The number of neighbors in KNN is referred to as a \emph{tuning
parameter}, also referred to as a \emph{hyperparameter}. We do not know
\emph{a priori} what value to use. It is therefore of interest to see
how the classifier performs on test data as we vary these parameters.
This can be achieved with a \texttt{for} loop, described in Section
2.3.8. Here we use a for loop to look at the accuracy of our classifier
in the group predicted to purchase insurance as we vary the number of
neighbors from 1 to 5:

    \begin{tcolorbox}[breakable, size=fbox, boxrule=1pt, pad at break*=1mm,colback=cellbackground, colframe=cellborder]
\prompt{In}{incolor}{60}{\boxspacing}
\begin{Verbatim}[commandchars=\\\{\}]
\PY{k}{for} \PY{n}{K} \PY{o+ow}{in} \PY{n+nb}{range}\PY{p}{(}\PY{l+m+mi}{1}\PY{p}{,}\PY{l+m+mi}{6}\PY{p}{)}\PY{p}{:}
    \PY{n}{knn} \PY{o}{=} \PY{n}{KNeighborsClassifier}\PY{p}{(}\PY{n}{n\PYZus{}neighbors}\PY{o}{=}\PY{n}{K}\PY{p}{)}
    \PY{n}{knn\PYZus{}pred} \PY{o}{=} \PY{n}{knn}\PY{o}{.}\PY{n}{fit}\PY{p}{(}\PY{n}{X\PYZus{}train}\PY{p}{,} \PY{n}{y\PYZus{}train}\PY{p}{)}\PY{o}{.}\PY{n}{predict}\PY{p}{(}\PY{n}{X\PYZus{}test}\PY{p}{)}
    \PY{n}{C} \PY{o}{=} \PY{n}{confusion\PYZus{}table}\PY{p}{(}\PY{n}{knn\PYZus{}pred}\PY{p}{,} \PY{n}{y\PYZus{}test}\PY{p}{)}
    \PY{n}{templ} \PY{o}{=} \PY{p}{(}\PY{l+s+s1}{\PYZsq{}}\PY{l+s+s1}{K=}\PY{l+s+si}{\PYZob{}0:d\PYZcb{}}\PY{l+s+s1}{: \PYZsh{} predicted to rent: }\PY{l+s+si}{\PYZob{}1:\PYZgt{}2\PYZcb{}}\PY{l+s+s1}{,}\PY{l+s+s1}{\PYZsq{}} \PY{o}{+}
            \PY{l+s+s1}{\PYZsq{}}\PY{l+s+s1}{  \PYZsh{} who did rent }\PY{l+s+si}{\PYZob{}2:d\PYZcb{}}\PY{l+s+s1}{, accuracy }\PY{l+s+si}{\PYZob{}3:.1\PYZpc{}\PYZcb{}}\PY{l+s+s1}{\PYZsq{}}\PY{p}{)}
    \PY{n}{pred} \PY{o}{=} \PY{n}{C}\PY{o}{.}\PY{n}{loc}\PY{p}{[}\PY{l+s+s1}{\PYZsq{}}\PY{l+s+s1}{Yes}\PY{l+s+s1}{\PYZsq{}}\PY{p}{]}\PY{o}{.}\PY{n}{sum}\PY{p}{(}\PY{p}{)}
    \PY{n}{did\PYZus{}rent} \PY{o}{=} \PY{n}{C}\PY{o}{.}\PY{n}{loc}\PY{p}{[}\PY{l+s+s1}{\PYZsq{}}\PY{l+s+s1}{Yes}\PY{l+s+s1}{\PYZsq{}}\PY{p}{,}\PY{l+s+s1}{\PYZsq{}}\PY{l+s+s1}{Yes}\PY{l+s+s1}{\PYZsq{}}\PY{p}{]}
    \PY{n+nb}{print}\PY{p}{(}\PY{n}{templ}\PY{o}{.}\PY{n}{format}\PY{p}{(}
          \PY{n}{K}\PY{p}{,}
          \PY{n}{pred}\PY{p}{,}
          \PY{n}{did\PYZus{}rent}\PY{p}{,}
          \PY{n}{did\PYZus{}rent} \PY{o}{/} \PY{n}{pred}\PY{p}{)}\PY{p}{)}
\end{Verbatim}
\end{tcolorbox}

    \begin{Verbatim}[commandchars=\\\{\}]
K=1: \# predicted to rent: 62,  \# who did rent 9, accuracy 14.5\%
K=2: \# predicted to rent:  6,  \# who did rent 1, accuracy 16.7\%
K=3: \# predicted to rent: 20,  \# who did rent 3, accuracy 15.0\%
K=4: \# predicted to rent:  4,  \# who did rent 0, accuracy 0.0\%
K=5: \# predicted to rent:  7,  \# who did rent 1, accuracy 14.3\%
    \end{Verbatim}

    We see some variability --- the numbers for \texttt{K=4} are very
different from the rest.

\hypertarget{comparison-to-logistic-regression}{%
\subsubsection{Comparison to Logistic
Regression}\label{comparison-to-logistic-regression}}

As a comparison, we can also fit a logistic regression model to the
data. This can also be done with \texttt{sklearn}, though by default it
fits something like the \emph{ridge regression} version of logistic
regression, which we introduce in Chapter 6. This can be modified by
appropriately setting the argument \texttt{C} below. Its default value
is 1 but by setting it to a very large number, the algorithm converges
to the same solution as the usual (unregularized) logistic regression
estimator discussed above.

Unlike the \texttt{statsmodels} package, \texttt{sklearn} focuses less
on inference and more on classification. Hence, the \texttt{summary}
methods seen in \texttt{statsmodels} and our simplified version seen
with \texttt{summarize} are not generally available for the classifiers
in \texttt{sklearn}.

    \begin{tcolorbox}[breakable, size=fbox, boxrule=1pt, pad at break*=1mm,colback=cellbackground, colframe=cellborder]
\prompt{In}{incolor}{61}{\boxspacing}
\begin{Verbatim}[commandchars=\\\{\}]
\PY{n}{logit} \PY{o}{=} \PY{n}{LogisticRegression}\PY{p}{(}\PY{n}{C}\PY{o}{=}\PY{l+m+mf}{1e10}\PY{p}{,} \PY{n}{solver}\PY{o}{=}\PY{l+s+s1}{\PYZsq{}}\PY{l+s+s1}{liblinear}\PY{l+s+s1}{\PYZsq{}}\PY{p}{)}
\PY{n}{logit}\PY{o}{.}\PY{n}{fit}\PY{p}{(}\PY{n}{X\PYZus{}train}\PY{p}{,} \PY{n}{y\PYZus{}train}\PY{p}{)}
\PY{n}{logit\PYZus{}pred} \PY{o}{=} \PY{n}{logit}\PY{o}{.}\PY{n}{predict\PYZus{}proba}\PY{p}{(}\PY{n}{X\PYZus{}test}\PY{p}{)}
\PY{n}{logit\PYZus{}labels} \PY{o}{=} \PY{n}{np}\PY{o}{.}\PY{n}{where}\PY{p}{(}\PY{n}{logit\PYZus{}pred}\PY{p}{[}\PY{p}{:}\PY{p}{,}\PY{l+m+mi}{1}\PY{p}{]} \PY{o}{\PYZgt{}} \PY{l+m+mf}{.5}\PY{p}{,} \PY{l+s+s1}{\PYZsq{}}\PY{l+s+s1}{Yes}\PY{l+s+s1}{\PYZsq{}}\PY{p}{,} \PY{l+s+s1}{\PYZsq{}}\PY{l+s+s1}{No}\PY{l+s+s1}{\PYZsq{}}\PY{p}{)}
\PY{n}{confusion\PYZus{}table}\PY{p}{(}\PY{n}{logit\PYZus{}labels}\PY{p}{,} \PY{n}{y\PYZus{}test}\PY{p}{)}
\end{Verbatim}
\end{tcolorbox}

            \begin{tcolorbox}[breakable, size=fbox, boxrule=.5pt, pad at break*=1mm, opacityfill=0]
\prompt{Out}{outcolor}{61}{\boxspacing}
\begin{Verbatim}[commandchars=\\\{\}]
Truth       No  Yes
Predicted
No         931   67
Yes          2    0
\end{Verbatim}
\end{tcolorbox}
        
    We used the argument
\texttt{solver=\textquotesingle{}liblinear\textquotesingle{}} above to
avoid a warning with the default solver which would indicate that the
algorithm does not converge.

If we use \(0.5\) as the predicted probability cut-off for the
classifier, then we have a problem: only two of the test observations
are predicted to purchase insurance. However, we are not required to use
a cut-off of \(0.5\). If we instead predict a purchase any time the
predicted probability of purchase exceeds \(0.25\), we get much better
results: we predict that 29 people will purchase insurance, and we are
correct for about 31\% of these people. This is almost five times better
than random guessing!

    \begin{tcolorbox}[breakable, size=fbox, boxrule=1pt, pad at break*=1mm,colback=cellbackground, colframe=cellborder]
\prompt{In}{incolor}{62}{\boxspacing}
\begin{Verbatim}[commandchars=\\\{\}]
\PY{n}{logit\PYZus{}labels} \PY{o}{=} \PY{n}{np}\PY{o}{.}\PY{n}{where}\PY{p}{(}\PY{n}{logit\PYZus{}pred}\PY{p}{[}\PY{p}{:}\PY{p}{,}\PY{l+m+mi}{1}\PY{p}{]}\PY{o}{\PYZgt{}}\PY{l+m+mf}{0.25}\PY{p}{,} \PY{l+s+s1}{\PYZsq{}}\PY{l+s+s1}{Yes}\PY{l+s+s1}{\PYZsq{}}\PY{p}{,} \PY{l+s+s1}{\PYZsq{}}\PY{l+s+s1}{No}\PY{l+s+s1}{\PYZsq{}}\PY{p}{)}
\PY{n}{confusion\PYZus{}table}\PY{p}{(}\PY{n}{logit\PYZus{}labels}\PY{p}{,} \PY{n}{y\PYZus{}test}\PY{p}{)}
\end{Verbatim}
\end{tcolorbox}

            \begin{tcolorbox}[breakable, size=fbox, boxrule=.5pt, pad at break*=1mm, opacityfill=0]
\prompt{Out}{outcolor}{62}{\boxspacing}
\begin{Verbatim}[commandchars=\\\{\}]
Truth       No  Yes
Predicted
No         913   58
Yes         20    9
\end{Verbatim}
\end{tcolorbox}
        
    \begin{tcolorbox}[breakable, size=fbox, boxrule=1pt, pad at break*=1mm,colback=cellbackground, colframe=cellborder]
\prompt{In}{incolor}{63}{\boxspacing}
\begin{Verbatim}[commandchars=\\\{\}]
\PY{l+m+mi}{9}\PY{o}{/}\PY{p}{(}\PY{l+m+mi}{20}\PY{o}{+}\PY{l+m+mi}{9}\PY{p}{)}
\end{Verbatim}
\end{tcolorbox}

            \begin{tcolorbox}[breakable, size=fbox, boxrule=.5pt, pad at break*=1mm, opacityfill=0]
\prompt{Out}{outcolor}{63}{\boxspacing}
\begin{Verbatim}[commandchars=\\\{\}]
0.3103448275862069
\end{Verbatim}
\end{tcolorbox}
        
    \hypertarget{linear-and-poisson-regression-on-the-bikeshare-data}{%
\subsection{Linear and Poisson Regression on the Bikeshare
Data}\label{linear-and-poisson-regression-on-the-bikeshare-data}}

Here we fit linear and Poisson regression models to the
\texttt{Bikeshare} data, as described in Section 4.6. The response
\texttt{bikers} measures the number of bike rentals per hour in
Washington, DC in the period 2010--2012.

    \begin{tcolorbox}[breakable, size=fbox, boxrule=1pt, pad at break*=1mm,colback=cellbackground, colframe=cellborder]
\prompt{In}{incolor}{64}{\boxspacing}
\begin{Verbatim}[commandchars=\\\{\}]
\PY{n}{Bike} \PY{o}{=} \PY{n}{load\PYZus{}data}\PY{p}{(}\PY{l+s+s1}{\PYZsq{}}\PY{l+s+s1}{Bikeshare}\PY{l+s+s1}{\PYZsq{}}\PY{p}{)}
\end{Verbatim}
\end{tcolorbox}

    Let's have a peek at the dimensions and names of the variables in this
dataframe.

    \begin{tcolorbox}[breakable, size=fbox, boxrule=1pt, pad at break*=1mm,colback=cellbackground, colframe=cellborder]
\prompt{In}{incolor}{65}{\boxspacing}
\begin{Verbatim}[commandchars=\\\{\}]
\PY{n}{Bike}\PY{o}{.}\PY{n}{shape}\PY{p}{,} \PY{n}{Bike}\PY{o}{.}\PY{n}{columns}
\end{Verbatim}
\end{tcolorbox}

            \begin{tcolorbox}[breakable, size=fbox, boxrule=.5pt, pad at break*=1mm, opacityfill=0]
\prompt{Out}{outcolor}{65}{\boxspacing}
\begin{Verbatim}[commandchars=\\\{\}]
((8645, 15),
 Index(['season', 'mnth', 'day', 'hr', 'holiday', 'weekday', 'workingday',
        'weathersit', 'temp', 'atemp', 'hum', 'windspeed', 'casual',
        'registered', 'bikers'],
       dtype='object'))
\end{Verbatim}
\end{tcolorbox}
        
    \hypertarget{linear-regression}{%
\subsubsection{Linear Regression}\label{linear-regression}}

We begin by fitting a linear regression model to the data.

    \begin{tcolorbox}[breakable, size=fbox, boxrule=1pt, pad at break*=1mm,colback=cellbackground, colframe=cellborder]
\prompt{In}{incolor}{66}{\boxspacing}
\begin{Verbatim}[commandchars=\\\{\}]
\PY{n}{X} \PY{o}{=} \PY{n}{MS}\PY{p}{(}\PY{p}{[}\PY{l+s+s1}{\PYZsq{}}\PY{l+s+s1}{mnth}\PY{l+s+s1}{\PYZsq{}}\PY{p}{,}
        \PY{l+s+s1}{\PYZsq{}}\PY{l+s+s1}{hr}\PY{l+s+s1}{\PYZsq{}}\PY{p}{,}
        \PY{l+s+s1}{\PYZsq{}}\PY{l+s+s1}{workingday}\PY{l+s+s1}{\PYZsq{}}\PY{p}{,}
        \PY{l+s+s1}{\PYZsq{}}\PY{l+s+s1}{temp}\PY{l+s+s1}{\PYZsq{}}\PY{p}{,}
        \PY{l+s+s1}{\PYZsq{}}\PY{l+s+s1}{weathersit}\PY{l+s+s1}{\PYZsq{}}\PY{p}{]}\PY{p}{)}\PY{o}{.}\PY{n}{fit\PYZus{}transform}\PY{p}{(}\PY{n}{Bike}\PY{p}{)}
\PY{n}{Y} \PY{o}{=} \PY{n}{Bike}\PY{p}{[}\PY{l+s+s1}{\PYZsq{}}\PY{l+s+s1}{bikers}\PY{l+s+s1}{\PYZsq{}}\PY{p}{]}
\PY{n}{M\PYZus{}lm} \PY{o}{=} \PY{n}{sm}\PY{o}{.}\PY{n}{OLS}\PY{p}{(}\PY{n}{Y}\PY{p}{,} \PY{n}{X}\PY{p}{)}\PY{o}{.}\PY{n}{fit}\PY{p}{(}\PY{p}{)}
\PY{n}{summarize}\PY{p}{(}\PY{n}{M\PYZus{}lm}\PY{p}{)}
\end{Verbatim}
\end{tcolorbox}

            \begin{tcolorbox}[breakable, size=fbox, boxrule=.5pt, pad at break*=1mm, opacityfill=0]
\prompt{Out}{outcolor}{66}{\boxspacing}
\begin{Verbatim}[commandchars=\\\{\}]
                                 coef  std err       t  P>|t|
intercept                    -68.6317    5.307 -12.932  0.000
mnth[Feb]                      6.8452    4.287   1.597  0.110
mnth[March]                   16.5514    4.301   3.848  0.000
mnth[April]                   41.4249    4.972   8.331  0.000
mnth[May]                     72.5571    5.641  12.862  0.000
mnth[June]                    67.8187    6.544  10.364  0.000
mnth[July]                    45.3245    7.081   6.401  0.000
mnth[Aug]                     53.2430    6.640   8.019  0.000
mnth[Sept]                    66.6783    5.925  11.254  0.000
mnth[Oct]                     75.8343    4.950  15.319  0.000
mnth[Nov]                     60.3100    4.610  13.083  0.000
mnth[Dec]                     46.4577    4.271  10.878  0.000
hr[1]                        -14.5793    5.699  -2.558  0.011
hr[2]                        -21.5791    5.733  -3.764  0.000
hr[3]                        -31.1408    5.778  -5.389  0.000
hr[4]                        -36.9075    5.802  -6.361  0.000
hr[5]                        -24.1355    5.737  -4.207  0.000
hr[6]                         20.5997    5.704   3.612  0.000
hr[7]                        120.0931    5.693  21.095  0.000
hr[8]                        223.6619    5.690  39.310  0.000
hr[9]                        120.5819    5.693  21.182  0.000
hr[10]                        83.8013    5.705  14.689  0.000
hr[11]                       105.4234    5.722  18.424  0.000
hr[12]                       137.2837    5.740  23.916  0.000
hr[13]                       136.0359    5.760  23.617  0.000
hr[14]                       126.6361    5.776  21.923  0.000
hr[15]                       132.0865    5.780  22.852  0.000
hr[16]                       178.5206    5.772  30.927  0.000
hr[17]                       296.2670    5.749  51.537  0.000
hr[18]                       269.4409    5.736  46.976  0.000
hr[19]                       186.2558    5.714  32.596  0.000
hr[20]                       125.5492    5.704  22.012  0.000
hr[21]                        87.5537    5.693  15.378  0.000
hr[22]                        59.1226    5.689  10.392  0.000
hr[23]                        26.8376    5.688   4.719  0.000
workingday                     1.2696    1.784   0.711  0.477
temp                         157.2094   10.261  15.321  0.000
weathersit[cloudy/misty]     -12.8903    1.964  -6.562  0.000
weathersit[heavy rain/snow] -109.7446   76.667  -1.431  0.152
weathersit[light rain/snow]  -66.4944    2.965 -22.425  0.000
\end{Verbatim}
\end{tcolorbox}
        
    There are 24 levels in \texttt{hr} and 40 rows in all. In
\texttt{M\_lm}, the first levels \texttt{hr{[}0{]}} and
\texttt{mnth{[}Jan{]}} are treated as the baseline values, and so no
coefficient estimates are provided for them: implicitly, their
coefficient estimates are zero, and all other levels are measured
relative to these baselines. For example, the Feb coefficient of
\(6.845\) signifies that, holding all other variables constant, there
are on average about 7 more riders in February than in January.
Similarly there are about 16.5 more riders in March than in January.

The results seen in Section 4.6.1 used a slightly different coding of
the variables \texttt{hr} and \texttt{mnth}, as follows:

    \begin{tcolorbox}[breakable, size=fbox, boxrule=1pt, pad at break*=1mm,colback=cellbackground, colframe=cellborder]
\prompt{In}{incolor}{67}{\boxspacing}
\begin{Verbatim}[commandchars=\\\{\}]
\PY{n}{hr\PYZus{}encode} \PY{o}{=} \PY{n}{contrast}\PY{p}{(}\PY{l+s+s1}{\PYZsq{}}\PY{l+s+s1}{hr}\PY{l+s+s1}{\PYZsq{}}\PY{p}{,} \PY{l+s+s1}{\PYZsq{}}\PY{l+s+s1}{sum}\PY{l+s+s1}{\PYZsq{}}\PY{p}{)}
\PY{n}{mnth\PYZus{}encode} \PY{o}{=} \PY{n}{contrast}\PY{p}{(}\PY{l+s+s1}{\PYZsq{}}\PY{l+s+s1}{mnth}\PY{l+s+s1}{\PYZsq{}}\PY{p}{,} \PY{l+s+s1}{\PYZsq{}}\PY{l+s+s1}{sum}\PY{l+s+s1}{\PYZsq{}}\PY{p}{)}
\end{Verbatim}
\end{tcolorbox}

    Refitting again:

    \begin{tcolorbox}[breakable, size=fbox, boxrule=1pt, pad at break*=1mm,colback=cellbackground, colframe=cellborder]
\prompt{In}{incolor}{68}{\boxspacing}
\begin{Verbatim}[commandchars=\\\{\}]
\PY{n}{X2} \PY{o}{=} \PY{n}{MS}\PY{p}{(}\PY{p}{[}\PY{n}{mnth\PYZus{}encode}\PY{p}{,}
         \PY{n}{hr\PYZus{}encode}\PY{p}{,}
        \PY{l+s+s1}{\PYZsq{}}\PY{l+s+s1}{workingday}\PY{l+s+s1}{\PYZsq{}}\PY{p}{,}
        \PY{l+s+s1}{\PYZsq{}}\PY{l+s+s1}{temp}\PY{l+s+s1}{\PYZsq{}}\PY{p}{,}
        \PY{l+s+s1}{\PYZsq{}}\PY{l+s+s1}{weathersit}\PY{l+s+s1}{\PYZsq{}}\PY{p}{]}\PY{p}{)}\PY{o}{.}\PY{n}{fit\PYZus{}transform}\PY{p}{(}\PY{n}{Bike}\PY{p}{)}
\PY{n}{M2\PYZus{}lm} \PY{o}{=} \PY{n}{sm}\PY{o}{.}\PY{n}{OLS}\PY{p}{(}\PY{n}{Y}\PY{p}{,} \PY{n}{X2}\PY{p}{)}\PY{o}{.}\PY{n}{fit}\PY{p}{(}\PY{p}{)}
\PY{n}{S2} \PY{o}{=} \PY{n}{summarize}\PY{p}{(}\PY{n}{M2\PYZus{}lm}\PY{p}{)}
\PY{n}{S2}
\end{Verbatim}
\end{tcolorbox}

            \begin{tcolorbox}[breakable, size=fbox, boxrule=.5pt, pad at break*=1mm, opacityfill=0]
\prompt{Out}{outcolor}{68}{\boxspacing}
\begin{Verbatim}[commandchars=\\\{\}]
                                 coef  std err       t  P>|t|
intercept                     73.5974    5.132  14.340  0.000
mnth[Jan]                    -46.0871    4.085 -11.281  0.000
mnth[Feb]                    -39.2419    3.539 -11.088  0.000
mnth[March]                  -29.5357    3.155  -9.361  0.000
mnth[April]                   -4.6622    2.741  -1.701  0.089
mnth[May]                     26.4700    2.851   9.285  0.000
mnth[June]                    21.7317    3.465   6.272  0.000
mnth[July]                    -0.7626    3.908  -0.195  0.845
mnth[Aug]                      7.1560    3.535   2.024  0.043
mnth[Sept]                    20.5912    3.046   6.761  0.000
mnth[Oct]                     29.7472    2.700  11.019  0.000
mnth[Nov]                     14.2229    2.860   4.972  0.000
hr[0]                        -96.1420    3.955 -24.307  0.000
hr[1]                       -110.7213    3.966 -27.916  0.000
hr[2]                       -117.7212    4.016 -29.310  0.000
hr[3]                       -127.2828    4.081 -31.191  0.000
hr[4]                       -133.0495    4.117 -32.319  0.000
hr[5]                       -120.2775    4.037 -29.794  0.000
hr[6]                        -75.5424    3.992 -18.925  0.000
hr[7]                         23.9511    3.969   6.035  0.000
hr[8]                        127.5199    3.950  32.284  0.000
hr[9]                         24.4399    3.936   6.209  0.000
hr[10]                       -12.3407    3.936  -3.135  0.002
hr[11]                         9.2814    3.945   2.353  0.019
hr[12]                        41.1417    3.957  10.397  0.000
hr[13]                        39.8939    3.975  10.036  0.000
hr[14]                        30.4940    3.991   7.641  0.000
hr[15]                        35.9445    3.995   8.998  0.000
hr[16]                        82.3786    3.988  20.655  0.000
hr[17]                       200.1249    3.964  50.488  0.000
hr[18]                       173.2989    3.956  43.806  0.000
hr[19]                        90.1138    3.940  22.872  0.000
hr[20]                        29.4071    3.936   7.471  0.000
hr[21]                        -8.5883    3.933  -2.184  0.029
hr[22]                       -37.0194    3.934  -9.409  0.000
workingday                     1.2696    1.784   0.711  0.477
temp                         157.2094   10.261  15.321  0.000
weathersit[cloudy/misty]     -12.8903    1.964  -6.562  0.000
weathersit[heavy rain/snow] -109.7446   76.667  -1.431  0.152
weathersit[light rain/snow]  -66.4944    2.965 -22.425  0.000
\end{Verbatim}
\end{tcolorbox}
        
    What is the difference between the two codings? In \texttt{M2\_lm}, a
coefficient estimate is reported for all but level \texttt{23} of
\texttt{hr} and level \texttt{Dec} of \texttt{mnth}. Importantly, in
\texttt{M2\_lm}, the (unreported) coefficient estimate for the last
level of \texttt{mnth} is not zero: instead, it equals the negative of
the sum of the coefficient estimates for all of the other levels.
Similarly, in \texttt{M2\_lm}, the coefficient estimate for the last
level of \texttt{hr} is the negative of the sum of the coefficient
estimates for all of the other levels. This means that the coefficients
of \texttt{hr} and \texttt{mnth} in \texttt{M2\_lm} will always sum to
zero, and can be interpreted as the difference from the mean level. For
example, the coefficient for January of \(-46.087\) indicates that,
holding all other variables constant, there are typically 46 fewer
riders in January relative to the yearly average.

It is important to realize that the choice of coding really does not
matter, provided that we interpret the model output correctly in light
of the coding used. For example, we see that the predictions from the
linear model are the same regardless of coding:

    \begin{tcolorbox}[breakable, size=fbox, boxrule=1pt, pad at break*=1mm,colback=cellbackground, colframe=cellborder]
\prompt{In}{incolor}{69}{\boxspacing}
\begin{Verbatim}[commandchars=\\\{\}]
\PY{n}{np}\PY{o}{.}\PY{n}{sum}\PY{p}{(}\PY{p}{(}\PY{n}{M\PYZus{}lm}\PY{o}{.}\PY{n}{fittedvalues} \PY{o}{\PYZhy{}} \PY{n}{M2\PYZus{}lm}\PY{o}{.}\PY{n}{fittedvalues}\PY{p}{)}\PY{o}{*}\PY{o}{*}\PY{l+m+mi}{2}\PY{p}{)}
\end{Verbatim}
\end{tcolorbox}

            \begin{tcolorbox}[breakable, size=fbox, boxrule=.5pt, pad at break*=1mm, opacityfill=0]
\prompt{Out}{outcolor}{69}{\boxspacing}
\begin{Verbatim}[commandchars=\\\{\}]
1.0334731385542263e-18
\end{Verbatim}
\end{tcolorbox}
        
    The sum of squared differences is zero. We can also see this using the
\texttt{np.allclose()} function:

    \begin{tcolorbox}[breakable, size=fbox, boxrule=1pt, pad at break*=1mm,colback=cellbackground, colframe=cellborder]
\prompt{In}{incolor}{70}{\boxspacing}
\begin{Verbatim}[commandchars=\\\{\}]
\PY{n}{np}\PY{o}{.}\PY{n}{allclose}\PY{p}{(}\PY{n}{M\PYZus{}lm}\PY{o}{.}\PY{n}{fittedvalues}\PY{p}{,} \PY{n}{M2\PYZus{}lm}\PY{o}{.}\PY{n}{fittedvalues}\PY{p}{)}
\end{Verbatim}
\end{tcolorbox}

            \begin{tcolorbox}[breakable, size=fbox, boxrule=.5pt, pad at break*=1mm, opacityfill=0]
\prompt{Out}{outcolor}{70}{\boxspacing}
\begin{Verbatim}[commandchars=\\\{\}]
True
\end{Verbatim}
\end{tcolorbox}
        
    To reproduce the left-hand side of Figure 4.13 we must first obtain the
coefficient estimates associated with \texttt{mnth}. The coefficients
for January through November can be obtained directly from the
\texttt{M2\_lm} object. The coefficient for December must be explicitly
computed as the negative sum of all the other months. We first extract
all the coefficients for month from the coefficients of \texttt{M2\_lm}.

    \begin{tcolorbox}[breakable, size=fbox, boxrule=1pt, pad at break*=1mm,colback=cellbackground, colframe=cellborder]
\prompt{In}{incolor}{71}{\boxspacing}
\begin{Verbatim}[commandchars=\\\{\}]
\PY{n}{coef\PYZus{}month} \PY{o}{=} \PY{n}{S2}\PY{p}{[}\PY{n}{S2}\PY{o}{.}\PY{n}{index}\PY{o}{.}\PY{n}{str}\PY{o}{.}\PY{n}{contains}\PY{p}{(}\PY{l+s+s1}{\PYZsq{}}\PY{l+s+s1}{mnth}\PY{l+s+s1}{\PYZsq{}}\PY{p}{)}\PY{p}{]}\PY{p}{[}\PY{l+s+s1}{\PYZsq{}}\PY{l+s+s1}{coef}\PY{l+s+s1}{\PYZsq{}}\PY{p}{]}
\PY{n}{coef\PYZus{}month}
\end{Verbatim}
\end{tcolorbox}

            \begin{tcolorbox}[breakable, size=fbox, boxrule=.5pt, pad at break*=1mm, opacityfill=0]
\prompt{Out}{outcolor}{71}{\boxspacing}
\begin{Verbatim}[commandchars=\\\{\}]
mnth[Jan]     -46.0871
mnth[Feb]     -39.2419
mnth[March]   -29.5357
mnth[April]    -4.6622
mnth[May]      26.4700
mnth[June]     21.7317
mnth[July]     -0.7626
mnth[Aug]       7.1560
mnth[Sept]     20.5912
mnth[Oct]      29.7472
mnth[Nov]      14.2229
Name: coef, dtype: float64
\end{Verbatim}
\end{tcolorbox}
        
    Next, we append \texttt{Dec} as the negative of the sum of all other
months.

    \begin{tcolorbox}[breakable, size=fbox, boxrule=1pt, pad at break*=1mm,colback=cellbackground, colframe=cellborder]
\prompt{In}{incolor}{72}{\boxspacing}
\begin{Verbatim}[commandchars=\\\{\}]
\PY{n}{months} \PY{o}{=} \PY{n}{Bike}\PY{p}{[}\PY{l+s+s1}{\PYZsq{}}\PY{l+s+s1}{mnth}\PY{l+s+s1}{\PYZsq{}}\PY{p}{]}\PY{o}{.}\PY{n}{dtype}\PY{o}{.}\PY{n}{categories}
\PY{n}{coef\PYZus{}month} \PY{o}{=} \PY{n}{pd}\PY{o}{.}\PY{n}{concat}\PY{p}{(}\PY{p}{[}
                       \PY{n}{coef\PYZus{}month}\PY{p}{,}
                       \PY{n}{pd}\PY{o}{.}\PY{n}{Series}\PY{p}{(}\PY{p}{[}\PY{o}{\PYZhy{}}\PY{n}{coef\PYZus{}month}\PY{o}{.}\PY{n}{sum}\PY{p}{(}\PY{p}{)}\PY{p}{]}\PY{p}{,}
                                  \PY{n}{index}\PY{o}{=}\PY{p}{[}\PY{l+s+s1}{\PYZsq{}}\PY{l+s+s1}{mnth[Dec]}\PY{l+s+s1}{\PYZsq{}}
                                 \PY{p}{]}\PY{p}{)}
                       \PY{p}{]}\PY{p}{)}
\PY{n}{coef\PYZus{}month}
\end{Verbatim}
\end{tcolorbox}

            \begin{tcolorbox}[breakable, size=fbox, boxrule=.5pt, pad at break*=1mm, opacityfill=0]
\prompt{Out}{outcolor}{72}{\boxspacing}
\begin{Verbatim}[commandchars=\\\{\}]
mnth[Jan]     -46.0871
mnth[Feb]     -39.2419
mnth[March]   -29.5357
mnth[April]    -4.6622
mnth[May]      26.4700
mnth[June]     21.7317
mnth[July]     -0.7626
mnth[Aug]       7.1560
mnth[Sept]     20.5912
mnth[Oct]      29.7472
mnth[Nov]      14.2229
mnth[Dec]       0.3705
dtype: float64
\end{Verbatim}
\end{tcolorbox}
        
    Finally, to make the plot neater, we'll just use the first letter of
each month, which is the \(6\)th entry of each of the labels in the
index.

    \begin{tcolorbox}[breakable, size=fbox, boxrule=1pt, pad at break*=1mm,colback=cellbackground, colframe=cellborder]
\prompt{In}{incolor}{73}{\boxspacing}
\begin{Verbatim}[commandchars=\\\{\}]
\PY{n}{fig\PYZus{}month}\PY{p}{,} \PY{n}{ax\PYZus{}month} \PY{o}{=} \PY{n}{subplots}\PY{p}{(}\PY{n}{figsize}\PY{o}{=}\PY{p}{(}\PY{l+m+mi}{8}\PY{p}{,}\PY{l+m+mi}{8}\PY{p}{)}\PY{p}{)}
\PY{n}{x\PYZus{}month} \PY{o}{=} \PY{n}{np}\PY{o}{.}\PY{n}{arange}\PY{p}{(}\PY{n}{coef\PYZus{}month}\PY{o}{.}\PY{n}{shape}\PY{p}{[}\PY{l+m+mi}{0}\PY{p}{]}\PY{p}{)}
\PY{n}{ax\PYZus{}month}\PY{o}{.}\PY{n}{plot}\PY{p}{(}\PY{n}{x\PYZus{}month}\PY{p}{,} \PY{n}{coef\PYZus{}month}\PY{p}{,} \PY{n}{marker}\PY{o}{=}\PY{l+s+s1}{\PYZsq{}}\PY{l+s+s1}{o}\PY{l+s+s1}{\PYZsq{}}\PY{p}{,} \PY{n}{ms}\PY{o}{=}\PY{l+m+mi}{10}\PY{p}{)}
\PY{n}{ax\PYZus{}month}\PY{o}{.}\PY{n}{set\PYZus{}xticks}\PY{p}{(}\PY{n}{x\PYZus{}month}\PY{p}{)}
\PY{n}{ax\PYZus{}month}\PY{o}{.}\PY{n}{set\PYZus{}xticklabels}\PY{p}{(}\PY{p}{[}\PY{n}{l}\PY{p}{[}\PY{l+m+mi}{5}\PY{p}{]} \PY{k}{for} \PY{n}{l} \PY{o+ow}{in} \PY{n}{coef\PYZus{}month}\PY{o}{.}\PY{n}{index}\PY{p}{]}\PY{p}{,} \PY{n}{fontsize}\PY{o}{=}\PY{l+m+mi}{20}\PY{p}{)}
\PY{n}{ax\PYZus{}month}\PY{o}{.}\PY{n}{set\PYZus{}xlabel}\PY{p}{(}\PY{l+s+s1}{\PYZsq{}}\PY{l+s+s1}{Month}\PY{l+s+s1}{\PYZsq{}}\PY{p}{,} \PY{n}{fontsize}\PY{o}{=}\PY{l+m+mi}{20}\PY{p}{)}
\PY{n}{ax\PYZus{}month}\PY{o}{.}\PY{n}{set\PYZus{}ylabel}\PY{p}{(}\PY{l+s+s1}{\PYZsq{}}\PY{l+s+s1}{Coefficient}\PY{l+s+s1}{\PYZsq{}}\PY{p}{,} \PY{n}{fontsize}\PY{o}{=}\PY{l+m+mi}{20}\PY{p}{)}\PY{p}{;}
\end{Verbatim}
\end{tcolorbox}

    \begin{center}
    \adjustimage{max size={0.9\linewidth}{0.9\paperheight}}{artifacts/latex/Ch04-classification-lab_files/artifacts/latex/Ch04-classification-lab_150_0.png}
    \end{center}
    { \hspace*{\fill} \\}
    
    Reproducing the right-hand plot in Figure 4.13 follows a similar
process.

    \begin{tcolorbox}[breakable, size=fbox, boxrule=1pt, pad at break*=1mm,colback=cellbackground, colframe=cellborder]
\prompt{In}{incolor}{74}{\boxspacing}
\begin{Verbatim}[commandchars=\\\{\}]
\PY{n}{coef\PYZus{}hr} \PY{o}{=} \PY{n}{S2}\PY{p}{[}\PY{n}{S2}\PY{o}{.}\PY{n}{index}\PY{o}{.}\PY{n}{str}\PY{o}{.}\PY{n}{contains}\PY{p}{(}\PY{l+s+s1}{\PYZsq{}}\PY{l+s+s1}{hr}\PY{l+s+s1}{\PYZsq{}}\PY{p}{)}\PY{p}{]}\PY{p}{[}\PY{l+s+s1}{\PYZsq{}}\PY{l+s+s1}{coef}\PY{l+s+s1}{\PYZsq{}}\PY{p}{]}
\PY{n}{coef\PYZus{}hr} \PY{o}{=} \PY{n}{coef\PYZus{}hr}\PY{o}{.}\PY{n}{reindex}\PY{p}{(}\PY{p}{[}\PY{l+s+s1}{\PYZsq{}}\PY{l+s+s1}{hr[}\PY{l+s+si}{\PYZob{}0\PYZcb{}}\PY{l+s+s1}{]}\PY{l+s+s1}{\PYZsq{}}\PY{o}{.}\PY{n}{format}\PY{p}{(}\PY{n}{h}\PY{p}{)} \PY{k}{for} \PY{n}{h} \PY{o+ow}{in} \PY{n+nb}{range}\PY{p}{(}\PY{l+m+mi}{23}\PY{p}{)}\PY{p}{]}\PY{p}{)}
\PY{n}{coef\PYZus{}hr} \PY{o}{=} \PY{n}{pd}\PY{o}{.}\PY{n}{concat}\PY{p}{(}\PY{p}{[}\PY{n}{coef\PYZus{}hr}\PY{p}{,}
                     \PY{n}{pd}\PY{o}{.}\PY{n}{Series}\PY{p}{(}\PY{p}{[}\PY{o}{\PYZhy{}}\PY{n}{coef\PYZus{}hr}\PY{o}{.}\PY{n}{sum}\PY{p}{(}\PY{p}{)}\PY{p}{]}\PY{p}{,} \PY{n}{index}\PY{o}{=}\PY{p}{[}\PY{l+s+s1}{\PYZsq{}}\PY{l+s+s1}{hr[23]}\PY{l+s+s1}{\PYZsq{}}\PY{p}{]}\PY{p}{)}
                    \PY{p}{]}\PY{p}{)}
\end{Verbatim}
\end{tcolorbox}

    We now make the hour plot.

    \begin{tcolorbox}[breakable, size=fbox, boxrule=1pt, pad at break*=1mm,colback=cellbackground, colframe=cellborder]
\prompt{In}{incolor}{75}{\boxspacing}
\begin{Verbatim}[commandchars=\\\{\}]
\PY{n}{fig\PYZus{}hr}\PY{p}{,} \PY{n}{ax\PYZus{}hr} \PY{o}{=} \PY{n}{subplots}\PY{p}{(}\PY{n}{figsize}\PY{o}{=}\PY{p}{(}\PY{l+m+mi}{8}\PY{p}{,}\PY{l+m+mi}{8}\PY{p}{)}\PY{p}{)}
\PY{n}{x\PYZus{}hr} \PY{o}{=} \PY{n}{np}\PY{o}{.}\PY{n}{arange}\PY{p}{(}\PY{n}{coef\PYZus{}hr}\PY{o}{.}\PY{n}{shape}\PY{p}{[}\PY{l+m+mi}{0}\PY{p}{]}\PY{p}{)}
\PY{n}{ax\PYZus{}hr}\PY{o}{.}\PY{n}{plot}\PY{p}{(}\PY{n}{x\PYZus{}hr}\PY{p}{,} \PY{n}{coef\PYZus{}hr}\PY{p}{,} \PY{n}{marker}\PY{o}{=}\PY{l+s+s1}{\PYZsq{}}\PY{l+s+s1}{o}\PY{l+s+s1}{\PYZsq{}}\PY{p}{,} \PY{n}{ms}\PY{o}{=}\PY{l+m+mi}{10}\PY{p}{)}
\PY{n}{ax\PYZus{}hr}\PY{o}{.}\PY{n}{set\PYZus{}xticks}\PY{p}{(}\PY{n}{x\PYZus{}hr}\PY{p}{[}\PY{p}{:}\PY{p}{:}\PY{l+m+mi}{2}\PY{p}{]}\PY{p}{)}
\PY{n}{ax\PYZus{}hr}\PY{o}{.}\PY{n}{set\PYZus{}xticklabels}\PY{p}{(}\PY{n+nb}{range}\PY{p}{(}\PY{l+m+mi}{24}\PY{p}{)}\PY{p}{[}\PY{p}{:}\PY{p}{:}\PY{l+m+mi}{2}\PY{p}{]}\PY{p}{,} \PY{n}{fontsize}\PY{o}{=}\PY{l+m+mi}{20}\PY{p}{)}
\PY{n}{ax\PYZus{}hr}\PY{o}{.}\PY{n}{set\PYZus{}xlabel}\PY{p}{(}\PY{l+s+s1}{\PYZsq{}}\PY{l+s+s1}{Hour}\PY{l+s+s1}{\PYZsq{}}\PY{p}{,} \PY{n}{fontsize}\PY{o}{=}\PY{l+m+mi}{20}\PY{p}{)}
\PY{n}{ax\PYZus{}hr}\PY{o}{.}\PY{n}{set\PYZus{}ylabel}\PY{p}{(}\PY{l+s+s1}{\PYZsq{}}\PY{l+s+s1}{Coefficient}\PY{l+s+s1}{\PYZsq{}}\PY{p}{,} \PY{n}{fontsize}\PY{o}{=}\PY{l+m+mi}{20}\PY{p}{)}\PY{p}{;}
\end{Verbatim}
\end{tcolorbox}

    \begin{center}
    \adjustimage{max size={0.9\linewidth}{0.9\paperheight}}{artifacts/latex/Ch04-classification-lab_files/artifacts/latex/Ch04-classification-lab_154_0.png}
    \end{center}
    { \hspace*{\fill} \\}
    
    \hypertarget{poisson-regression}{%
\subsubsection{Poisson Regression}\label{poisson-regression}}

Now we fit instead a Poisson regression model to the \texttt{Bikeshare}
data. Very little changes, except that we now use the function
\texttt{sm.GLM()} with the Poisson family specified:

    \begin{tcolorbox}[breakable, size=fbox, boxrule=1pt, pad at break*=1mm,colback=cellbackground, colframe=cellborder]
\prompt{In}{incolor}{76}{\boxspacing}
\begin{Verbatim}[commandchars=\\\{\}]
\PY{n}{M\PYZus{}pois} \PY{o}{=} \PY{n}{sm}\PY{o}{.}\PY{n}{GLM}\PY{p}{(}\PY{n}{Y}\PY{p}{,} \PY{n}{X2}\PY{p}{,} \PY{n}{family}\PY{o}{=}\PY{n}{sm}\PY{o}{.}\PY{n}{families}\PY{o}{.}\PY{n}{Poisson}\PY{p}{(}\PY{p}{)}\PY{p}{)}\PY{o}{.}\PY{n}{fit}\PY{p}{(}\PY{p}{)}
\end{Verbatim}
\end{tcolorbox}

    We can plot the coefficients associated with \texttt{mnth} and
\texttt{hr}, in order to reproduce Figure 4.15. We first complete these
coefficients as before.

    \begin{tcolorbox}[breakable, size=fbox, boxrule=1pt, pad at break*=1mm,colback=cellbackground, colframe=cellborder]
\prompt{In}{incolor}{77}{\boxspacing}
\begin{Verbatim}[commandchars=\\\{\}]
\PY{n}{S\PYZus{}pois} \PY{o}{=} \PY{n}{summarize}\PY{p}{(}\PY{n}{M\PYZus{}pois}\PY{p}{)}
\PY{n}{coef\PYZus{}month} \PY{o}{=} \PY{n}{S\PYZus{}pois}\PY{p}{[}\PY{n}{S\PYZus{}pois}\PY{o}{.}\PY{n}{index}\PY{o}{.}\PY{n}{str}\PY{o}{.}\PY{n}{contains}\PY{p}{(}\PY{l+s+s1}{\PYZsq{}}\PY{l+s+s1}{mnth}\PY{l+s+s1}{\PYZsq{}}\PY{p}{)}\PY{p}{]}\PY{p}{[}\PY{l+s+s1}{\PYZsq{}}\PY{l+s+s1}{coef}\PY{l+s+s1}{\PYZsq{}}\PY{p}{]}
\PY{n}{coef\PYZus{}month} \PY{o}{=} \PY{n}{pd}\PY{o}{.}\PY{n}{concat}\PY{p}{(}\PY{p}{[}\PY{n}{coef\PYZus{}month}\PY{p}{,}
                        \PY{n}{pd}\PY{o}{.}\PY{n}{Series}\PY{p}{(}\PY{p}{[}\PY{o}{\PYZhy{}}\PY{n}{coef\PYZus{}month}\PY{o}{.}\PY{n}{sum}\PY{p}{(}\PY{p}{)}\PY{p}{]}\PY{p}{,}
                                   \PY{n}{index}\PY{o}{=}\PY{p}{[}\PY{l+s+s1}{\PYZsq{}}\PY{l+s+s1}{mnth[Dec]}\PY{l+s+s1}{\PYZsq{}}\PY{p}{]}\PY{p}{)}\PY{p}{]}\PY{p}{)}
\PY{n}{coef\PYZus{}hr} \PY{o}{=} \PY{n}{S\PYZus{}pois}\PY{p}{[}\PY{n}{S\PYZus{}pois}\PY{o}{.}\PY{n}{index}\PY{o}{.}\PY{n}{str}\PY{o}{.}\PY{n}{contains}\PY{p}{(}\PY{l+s+s1}{\PYZsq{}}\PY{l+s+s1}{hr}\PY{l+s+s1}{\PYZsq{}}\PY{p}{)}\PY{p}{]}\PY{p}{[}\PY{l+s+s1}{\PYZsq{}}\PY{l+s+s1}{coef}\PY{l+s+s1}{\PYZsq{}}\PY{p}{]}
\PY{n}{coef\PYZus{}hr} \PY{o}{=} \PY{n}{pd}\PY{o}{.}\PY{n}{concat}\PY{p}{(}\PY{p}{[}\PY{n}{coef\PYZus{}hr}\PY{p}{,}
                     \PY{n}{pd}\PY{o}{.}\PY{n}{Series}\PY{p}{(}\PY{p}{[}\PY{o}{\PYZhy{}}\PY{n}{coef\PYZus{}hr}\PY{o}{.}\PY{n}{sum}\PY{p}{(}\PY{p}{)}\PY{p}{]}\PY{p}{,}
                     \PY{n}{index}\PY{o}{=}\PY{p}{[}\PY{l+s+s1}{\PYZsq{}}\PY{l+s+s1}{hr[23]}\PY{l+s+s1}{\PYZsq{}}\PY{p}{]}\PY{p}{)}\PY{p}{]}\PY{p}{)}
\end{Verbatim}
\end{tcolorbox}

    The plotting is as before.

    \begin{tcolorbox}[breakable, size=fbox, boxrule=1pt, pad at break*=1mm,colback=cellbackground, colframe=cellborder]
\prompt{In}{incolor}{78}{\boxspacing}
\begin{Verbatim}[commandchars=\\\{\}]
\PY{n}{fig\PYZus{}pois}\PY{p}{,} \PY{p}{(}\PY{n}{ax\PYZus{}month}\PY{p}{,} \PY{n}{ax\PYZus{}hr}\PY{p}{)} \PY{o}{=} \PY{n}{subplots}\PY{p}{(}\PY{l+m+mi}{1}\PY{p}{,} \PY{l+m+mi}{2}\PY{p}{,} \PY{n}{figsize}\PY{o}{=}\PY{p}{(}\PY{l+m+mi}{16}\PY{p}{,}\PY{l+m+mi}{8}\PY{p}{)}\PY{p}{)}
\PY{n}{ax\PYZus{}month}\PY{o}{.}\PY{n}{plot}\PY{p}{(}\PY{n}{x\PYZus{}month}\PY{p}{,} \PY{n}{coef\PYZus{}month}\PY{p}{,} \PY{n}{marker}\PY{o}{=}\PY{l+s+s1}{\PYZsq{}}\PY{l+s+s1}{o}\PY{l+s+s1}{\PYZsq{}}\PY{p}{,} \PY{n}{ms}\PY{o}{=}\PY{l+m+mi}{10}\PY{p}{)}
\PY{n}{ax\PYZus{}month}\PY{o}{.}\PY{n}{set\PYZus{}xticks}\PY{p}{(}\PY{n}{x\PYZus{}month}\PY{p}{)}
\PY{n}{ax\PYZus{}month}\PY{o}{.}\PY{n}{set\PYZus{}xticklabels}\PY{p}{(}\PY{p}{[}\PY{n}{l}\PY{p}{[}\PY{l+m+mi}{5}\PY{p}{]} \PY{k}{for} \PY{n}{l} \PY{o+ow}{in} \PY{n}{coef\PYZus{}month}\PY{o}{.}\PY{n}{index}\PY{p}{]}\PY{p}{,} \PY{n}{fontsize}\PY{o}{=}\PY{l+m+mi}{20}\PY{p}{)}
\PY{n}{ax\PYZus{}month}\PY{o}{.}\PY{n}{set\PYZus{}xlabel}\PY{p}{(}\PY{l+s+s1}{\PYZsq{}}\PY{l+s+s1}{Month}\PY{l+s+s1}{\PYZsq{}}\PY{p}{,} \PY{n}{fontsize}\PY{o}{=}\PY{l+m+mi}{20}\PY{p}{)}
\PY{n}{ax\PYZus{}month}\PY{o}{.}\PY{n}{set\PYZus{}ylabel}\PY{p}{(}\PY{l+s+s1}{\PYZsq{}}\PY{l+s+s1}{Coefficient}\PY{l+s+s1}{\PYZsq{}}\PY{p}{,} \PY{n}{fontsize}\PY{o}{=}\PY{l+m+mi}{20}\PY{p}{)}
\PY{n}{ax\PYZus{}hr}\PY{o}{.}\PY{n}{plot}\PY{p}{(}\PY{n}{x\PYZus{}hr}\PY{p}{,} \PY{n}{coef\PYZus{}hr}\PY{p}{,} \PY{n}{marker}\PY{o}{=}\PY{l+s+s1}{\PYZsq{}}\PY{l+s+s1}{o}\PY{l+s+s1}{\PYZsq{}}\PY{p}{,} \PY{n}{ms}\PY{o}{=}\PY{l+m+mi}{10}\PY{p}{)}
\PY{n}{ax\PYZus{}hr}\PY{o}{.}\PY{n}{set\PYZus{}xticklabels}\PY{p}{(}\PY{n+nb}{range}\PY{p}{(}\PY{l+m+mi}{24}\PY{p}{)}\PY{p}{[}\PY{p}{:}\PY{p}{:}\PY{l+m+mi}{2}\PY{p}{]}\PY{p}{,} \PY{n}{fontsize}\PY{o}{=}\PY{l+m+mi}{20}\PY{p}{)}
\PY{n}{ax\PYZus{}hr}\PY{o}{.}\PY{n}{set\PYZus{}xlabel}\PY{p}{(}\PY{l+s+s1}{\PYZsq{}}\PY{l+s+s1}{Hour}\PY{l+s+s1}{\PYZsq{}}\PY{p}{,} \PY{n}{fontsize}\PY{o}{=}\PY{l+m+mi}{20}\PY{p}{)}
\PY{n}{ax\PYZus{}hr}\PY{o}{.}\PY{n}{set\PYZus{}ylabel}\PY{p}{(}\PY{l+s+s1}{\PYZsq{}}\PY{l+s+s1}{Coefficient}\PY{l+s+s1}{\PYZsq{}}\PY{p}{,} \PY{n}{fontsize}\PY{o}{=}\PY{l+m+mi}{20}\PY{p}{)}\PY{p}{;}
\end{Verbatim}
\end{tcolorbox}

    \begin{Verbatim}[commandchars=\\\{\}]
/var/folders/16/8y65\_zv174qgdp4ktlmpv12h0000gq/T/ipykernel\_70549/3779905754.py:8
: UserWarning: set\_ticklabels() should only be used with a fixed number of
ticks, i.e. after set\_ticks() or using a FixedLocator.
  ax\_hr.set\_xticklabels(range(24)[::2], fontsize=20)
    \end{Verbatim}

    \begin{center}
    \adjustimage{max size={0.9\linewidth}{0.9\paperheight}}{artifacts/latex/Ch04-classification-lab_files/artifacts/latex/Ch04-classification-lab_160_1.png}
    \end{center}
    { \hspace*{\fill} \\}
    
    We compare the fitted values of the two models. The fitted values are
stored in the \texttt{fittedvalues} attribute returned by the
\texttt{fit()} method for both the linear regression and the Poisson
fits. The linear predictors are stored as the attribute
\texttt{lin\_pred}.

    \begin{tcolorbox}[breakable, size=fbox, boxrule=1pt, pad at break*=1mm,colback=cellbackground, colframe=cellborder]
\prompt{In}{incolor}{79}{\boxspacing}
\begin{Verbatim}[commandchars=\\\{\}]
\PY{n}{fig}\PY{p}{,} \PY{n}{ax} \PY{o}{=} \PY{n}{subplots}\PY{p}{(}\PY{n}{figsize}\PY{o}{=}\PY{p}{(}\PY{l+m+mi}{8}\PY{p}{,} \PY{l+m+mi}{8}\PY{p}{)}\PY{p}{)}
\PY{n}{ax}\PY{o}{.}\PY{n}{scatter}\PY{p}{(}\PY{n}{M2\PYZus{}lm}\PY{o}{.}\PY{n}{fittedvalues}\PY{p}{,}
           \PY{n}{M\PYZus{}pois}\PY{o}{.}\PY{n}{fittedvalues}\PY{p}{,}
           \PY{n}{s}\PY{o}{=}\PY{l+m+mi}{20}\PY{p}{)}
\PY{n}{ax}\PY{o}{.}\PY{n}{set\PYZus{}xlabel}\PY{p}{(}\PY{l+s+s1}{\PYZsq{}}\PY{l+s+s1}{Linear Regression Fit}\PY{l+s+s1}{\PYZsq{}}\PY{p}{,} \PY{n}{fontsize}\PY{o}{=}\PY{l+m+mi}{20}\PY{p}{)}
\PY{n}{ax}\PY{o}{.}\PY{n}{set\PYZus{}ylabel}\PY{p}{(}\PY{l+s+s1}{\PYZsq{}}\PY{l+s+s1}{Poisson Regression Fit}\PY{l+s+s1}{\PYZsq{}}\PY{p}{,} \PY{n}{fontsize}\PY{o}{=}\PY{l+m+mi}{20}\PY{p}{)}
\PY{n}{ax}\PY{o}{.}\PY{n}{axline}\PY{p}{(}\PY{p}{[}\PY{l+m+mi}{0}\PY{p}{,}\PY{l+m+mi}{0}\PY{p}{]}\PY{p}{,} \PY{n}{c}\PY{o}{=}\PY{l+s+s1}{\PYZsq{}}\PY{l+s+s1}{black}\PY{l+s+s1}{\PYZsq{}}\PY{p}{,} \PY{n}{linewidth}\PY{o}{=}\PY{l+m+mi}{3}\PY{p}{,}
          \PY{n}{linestyle}\PY{o}{=}\PY{l+s+s1}{\PYZsq{}}\PY{l+s+s1}{\PYZhy{}\PYZhy{}}\PY{l+s+s1}{\PYZsq{}}\PY{p}{,} \PY{n}{slope}\PY{o}{=}\PY{l+m+mi}{1}\PY{p}{)}\PY{p}{;}
\end{Verbatim}
\end{tcolorbox}

    \begin{center}
    \adjustimage{max size={0.9\linewidth}{0.9\paperheight}}{artifacts/latex/Ch04-classification-lab_files/artifacts/latex/Ch04-classification-lab_162_0.png}
    \end{center}
    { \hspace*{\fill} \\}
    
    The predictions from the Poisson regression model are correlated with
those from the linear model; however, the former are non-negative. As a
result the Poisson regression predictions tend to be larger than those
from the linear model for either very low or very high levels of
ridership.

In this section, we fit Poisson regression models using the
\texttt{sm.GLM()} function with the argument
\texttt{family=sm.families.Poisson()}. Earlier in this lab we used the
\texttt{sm.GLM()} function with \texttt{family=sm.families.Binomial()}
to perform logistic regression. Other choices for the \texttt{family}
argument can be used to fit other types of GLMs. For instance,
\texttt{family=sm.families.Gamma()} fits a Gamma regression model.


    % Add a bibliography block to the postdoc
    
    
    
\end{document}
